\chapter{Caazapá}\label{dept:Caazapá}
\mapaparaguai{6}{Caazapá}
\vfill\clearpage
\section{Caazapá}

base agropecuária de baixa densidade e serviços concentrados em Caazapa
e San Juan Nepomuceno.

\subsection{Ficha climática de Caazapa (capital
departamental)}

\textbf{Irradiância solar global --- (kWh/\allowbreak m²/\allowbreak dia):}


\begin{center}
\resizebox{\linewidth}{!}{%
{\renewcommand{\arraystretch}{1.2}%
\begin{tabular}{@{}c|c|c|c|c|c|c|c|c|c|c|c@{}}
\toprule
Jan & Fev & Mar & Abr & Mai & Jun & Jul & Ago & Set & Out & Nov & Dez \\
\midrule
6.71 & 6.20 & 5.39 & 4.34 & 3.26 & 2.77 & 3.10 & 3.79 & 4.42 & 5.34 &
6.39 & 6.76 \\
\bottomrule
\end{tabular}}%
}
\end{center}


Média anual: 4.87 kWh/\allowbreak m²/\allowbreak dia. Inclinação recomendada: 26° Norte (anual);
36° inverno; 16° verão.

\textbf{Precipitação --- (mm/\allowbreak mês):}


\begin{center}
\resizebox{\linewidth}{!}{%
{\renewcommand{\arraystretch}{1.2}%
\begin{tabular}{@{}c|c|c|c|c|c|c|c|c|c|c|c@{}}
\toprule
Jan & Fev & Mar & Abr & Mai & Jun & Jul & Ago & Set & Out & Nov & Dez \\
\midrule
135 & 129 & 133 & 159 & 166 & 94 & 78 & 61 & 93 & 189 & 188 & 164 \\
\bottomrule
\end{tabular}}%
}
\end{center}


Média anual: \textasciitilde1.587 mm/\allowbreak ano. Estação chuvosa Out-Mai; seca
Jun-Set.
\clearpage
\subsection*{Dados Departamentais}

\subsection*{Segurança Pública}

\begin{tabularx}{\linewidth}{lX}
\toprule
Indicador & Valor \\
\midrule
Taxa de homicídios (por 100k hab) & ~4--6 (estimativa, próxima à média nacional) \\
Crimes registrados (total/ano) & --- dados não desagregados publicamente \\
Número de comisarías & ~8 (estimativa departamental) \\
Índice segurança relativo & médio \\
\bottomrule
\end{tabularx}

Caazapá é um dos departamentos menores e menos populosos do Paraguai, com perfil de segurança relativamente moderado. O Ministerio Público registra que o departamento responde por \textbf{aproximadamente 3\% das causas de homicídio doloso} nacional (2019--2024) --- proporção ligeiramente acima do esperado para sua participação populacional (~2,5\%).

O departamento não possui fronteira internacional e está localizado na região sul do Paraguai. Não é identificado como rota principal de narcotráfico, embora a porosidade do território rural permita algum fluxo de drogas. A maioria dos crimes violentos ocorre em contexto interpessoal, seguindo o padrão nacional.

A infraestrutura policial é limitada dado o tamanho do território e a dispersão da população. Caazapá (cidade capital) é pequena com serviços básicos disponíveis. A pobreza estrutural do departamento --- um dos mais pobres do Paraguai --- contribui para algumas formas de criminalidade.

\subsection*{IDH e Desenvolvimento Humano}

\begin{tabularx}{\linewidth}{lX}
\toprule
Indicador & Valor \\
\midrule
IDH & 0,648 \\
Classificação IDH & médio \\
Ranking nacional (1=melhor) & 18/18 \\
Esperança de vida & 69,5 anos \\
Anos de escolaridade média & 6,4 anos \\
RNB per capita (USD PPA) & 5.200 \\
\bottomrule
\end{tabularx}

> Nota: IDH de 2020 conforme relatório PNUD "Desarrollo Humano en el Paraguay 2001-2020". Caazapá é explicitamente citado como o departamento de menor IDH do Paraguai, com a brecha de 42 anos em relação a Asunción (ABC Color, fevereiro 2023). Fonte: PNUD Paraguay, ABC Color.

\begin{tabularx}{\linewidth}{lX}
\toprule
Indicador & Valor \\
\midrule
Taxa de pobreza total (\%) & 34,7\% \\
Taxa de pobreza extrema (\%) & 11,3\% \\
Índice de Gini & 0,472 \\
\bottomrule
\end{tabularx}

> Fonte: INE --- EPHC 2023. Caazapá possui consistentemente a maior taxa de pobreza percentual entre os departamentos com dados históricos. Referência citada: ABC Color (pobreza 36,8\% em período anterior, estável no intervalo 34--37\%).

\begin{tabularx}{\linewidth}{lX}
\toprule
Indicador & Valor \\
\midrule
Acesso a água potável (\%) & 60,1\% \\
Acesso a saneamento (\%) & 57,4\% \\
\bottomrule
\end{tabularx}

> Estimativa baseada em dados do Censo 2022 e perfil predominantemente rural. Caazapá tem 80\%+ da população em zona rural, com cobertura de serviços básicos muito abaixo da média nacional.

Caazapá ocupa o último lugar no IDH departamental do Paraguai (0,648) e apresenta os maiores desafios em termos de desenvolvimento. A brecha em relação a Asunción é tão grande que levaria 42 anos de crescimento contínuo para igualar os indicadores da capital, conforme estudo do PNUD. Para imigrantes brasileiros, o departamento oferece terras muito baratas para agricultura, mas a infraestrutura é precária: hospitais com recursos limitados, estradas em condições irregulares, acesso restrito a água tratada (60,1\%) e saneamento (57,4\%). O custo de vida é extremamente baixo, mas em contrapartida os serviços são escassos. Recomendado apenas para projetos de agricultura extensiva de médio/longo prazo por quem está disposto a operar em ambiente de infraestrutura limitada. Potencial agrícola não explorado com possibilidade de valorização futura.

\subsection*{Sistema Prisional}

{\small
{\small
\begin{tabularx}{\linewidth}{lXXXX}
\toprule
Nome & Município & Capacidade & População & Superlotação \\
\midrule
Penitenciaría Regional de Caazapá & Caazapá & 100 & ~185 & ~185\% \\
\bottomrule
\end{tabularx}
}
}

\begin{tabularx}{\linewidth}{lX}
\toprule
Indicador & Valor \\
\midrule
Total de estabelecimentos & 1 \\
Capacidade total (vagas) & ~100 \\
População carcerária total & ~185 \\
Taxa de superlotação & ~185\% \\
\% presos provisórios & ~68\% \\
\bottomrule
\end{tabularx}

O departamento de Caazapá possui apenas uma penitenciaría regional, localizada na capital departamental Caazapá. É um dos menores estabelecimentos penais do Paraguai em termos de capacidade, mas ainda assim apresenta superlotação significativa.

Caazapá é um dos departamentos mais pobres do Paraguai, com economia baseada principalmente na agricultura familiar e pecuária extensiva. O perfil da criminalidade local é predominantemente relacionado a crimes contra a propriedade e conflitos fundiários --- herança de tensões históricas pela posse de terras.

O MNP relatou em inspeções que o estabelecimento de Caazapá carece de infraestrutura mínima adequada: celas superlotadas, ventilação insuficiente e serviços médicos precários. Detentos de municípios mais distantes (Yuty, General Higinio Morínigo) podem levar horas de deslocamento até o presídio para visitas familiares.

\textbf{Para imigrantes:} Caazapá oferece baixo custo de vida e terras disponíveis para agricultura, mas tem infraestrutura urbana limitada e menor acesso a serviços. O sistema prisional não representa risco direto para a maioria das atividades lícitas no departamento. Recomendado para quem busca vida rural tranquila.

\subsection*{Recursos Hídricos Subterrâneos}

\begin{tabularx}{\linewidth}{lXXX}
\toprule
Aquífero & Cobertura no Dept. & Profundidade Média & Vazão Típica \\
\midrule
Guarani / Misiones (confinado) & ~75\% & 80--300 m & 20--80 m$^3$/h \\
Cristalino (fraturado) & ~25\% (porção oeste) & 30--80 m & 2--8 m$^3$/h \\
\bottomrule
\end{tabularx}

\begin{tabularx}{\linewidth}{lX}
\toprule
Parâmetro & Valor \\
\midrule
Profundidade média --- Guarani (m) & 80--300 \\
Profundidade máxima recomendada (m) & 400 \\
Vazão típica (m$^3$/h) & 15--60 \\
Custo estimado perfuração (USD) & 3.000 -- 10.000 \\
Qualidade da água & boa a excelente \\
pH médio & 6,5--7,5 \\
Outorga necessária & sim \\
Órgão regulador & MADES --- DGPCRH \\
\bottomrule
\end{tabularx}

Caazapá está situado no coração da Região Oriental do Paraguai, com ampla cobertura do Aquífero Guarani/Misiones. O departamento apresenta afloramentos de arenito Misiones em vários pontos, o que facilita a recarga e a captação em profundidades moderadas. A água é de boa a excelente qualidade, com baixa mineralização, pH neutro e ausência de contaminantes naturais. O uso principal é abastecimento doméstico rural e dessedentação animal, com crescente uso para irrigação de soja e trigo. As cidades de Caazapá (capital) e Yuty são atendidas parcialmente por redes do SENASA.

Caazapá é um dos departamentos mais rurais do Paraguai, com baixa densidade populacional e terras relativamente baratas. Para propriedades rurais, o Aquífero Guarani é amplamente acessível (80--200 m), com custo de perfuração estimado em USD 3.000 a 8.000. A rede pública é limitada fora dos centros urbanos, tornando o poço artesiano indispensável. Excelente custo-benefício para uso doméstico e pequena irrigação. Licença do MADES é obrigatória antes da perfuração.

\subsection*{Aptidão do Solo}

\begin{tabularx}{\linewidth}{lXX}
\toprule
Tipo de Solo & \% do Território & Características \\
\midrule
Oxissolos (Latossolos Roxos) & 35\% & Solos derivados de basalto, profundos e bem drenados, alta aptidão agrícola \\
Ultissolos (Argissolos) & 35\% & Solos podzólicos de fertilidade média; ocorrem em ampla faixa central \\
Entissolos arenosos & 20\% & Solos arenosos de baixa fertilidade natural, suscetíveis à erosão; norte e leste \\
Inceptissolos / Aluviais & 10\% & Planícies fluviais do Rio Tebicuary e afluentes \\
\bottomrule
\end{tabularx}

\begin{tabularx}{\linewidth}{lXX}
\toprule
Classe & \% Território & Uso Recomendado \\
\midrule
Lavoura intensiva & 30\% & Soja, milho, trigo (solos basálticos) \\
Lavoura extensiva & 30\% & Mandioca, amendoim, algodão, pastagem melhorada \\
Pastagem natural & 20\% & Pecuária extensiva \\
Floresta / reserva & 15\% & Reservas nativas, matas ciliares \\
Sem aptidão agrícola & 5\% & Áreas urbanas e degradadas \\
\bottomrule
\end{tabularx}

\begin{tabularx}{\linewidth}{lX}
\toprule
Parâmetro & Valor \\
\midrule
pH médio & 4,8 -- 5,8 \\
Fertilidade natural & baixa a média \\
Teor de matéria orgânica & baixo a médio \\
Necessidade de calcário & sim \\
\bottomrule
\end{tabularx}

Caazapá está emergindo como uma das fronteiras agrícolas mais atrativas do Paraguai. As zonas com Oxissolos são comparáveis em qualidade aos Latossolos Roxos do sul do Brasil. O preço da terra ainda é inferior ao de Alto Paraná e Itapúa, com crescente interesse de produtores brasileiros. Pesquisa recente (MarketData, 2024) destaca Caazapá entre as três zonas mais produtivas para agricultura do país. A localização no corredor entre Encarnación e Villarrica facilita o escoamento. Requer investimento inicial em correção de solo.

\subsection*{Terra Rural}

\begin{tabularx}{\linewidth}{lXXX}
\toprule
Tipo & Preço (USD/ha) & Faixa & Tendência \\
\midrule
Agrícola alta produtividade & 5.500 & 3.500 -- 8.000 & $\uparrow$ \\
Pastagem & 2.500 & 1.800 -- 3.500 & $\uparrow$ \\
Mata / reserva & 1.000 & 600 -- 1.800 & $\rightarrow$ \\
\bottomrule
\end{tabularx}

\begin{tabularx}{\linewidth}{lX}
\toprule
Parâmetro & Valor \\
\midrule
Valorização anual média (\%) & 6--8\% \\
Lote mínimo rural (ha) & 20 \\
Imposto territorial rural (\%) & 1\% sobre valor fiscal \\
Liquidez do mercado & média \\
\bottomrule
\end{tabularx}

\begin{tabularx}{\linewidth}{lXXX}
\toprule
Indicador & Caazapá (PY) & Paraná (BR) & Rio Grande do Sul (BR) \\
\midrule
Terra agrícola (USD/ha) & 3.500--8.000 & 15.000--30.000 & 12.000--22.000 \\
Terra pastagem (USD/ha) & 1.800--3.500 & 6.000--12.000 & 5.000--10.000 \\
Custo relativo & ~30\% do PR & referência PR & referência RS \\
\bottomrule
\end{tabularx}

Caazapá é um departamento de fronteira agrícola em expansão, menos desenvolvido que Itapúa ou Alto Paraná, mas com solos férteis e preços ainda competitivos. Boa opção para quem busca maior área por menor investimento.

\subsection*{Mercado Imobiliário}

\textbf{Capital:} Caazapá

\begin{tabularx}{\linewidth}{lXX}
\toprule
Tipo & Preço (USD/m$^2$) & Aluguel (USD/mês) \\
\midrule
Apartamento 2 quartos & 380 -- 600 & 150 -- 280 \\
Casa 3 quartos & 320 -- 520 & 130 -- 250 \\
\bottomrule
\end{tabularx}

\begin{tabularx}{\linewidth}{lX}
\toprule
Parâmetro & Valor \\
\midrule
Valorização anual (\%) & 3--5\% \\
Disponibilidade de financiamento & limitada (cooperativas locais) \\
Bairros mais valorizados & Centro, Barrio San Roque, proximidades da Plaza Principal \\
\bottomrule
\end{tabularx}

\textbf{Vale a pena comprar ou alugar?}
Caazapá é uma das cidades mais tranquilas e menos populosas do Paraguai oriental (menos de 25.000 hab. no município). O mercado imobiliário é muito pequeno e informal. Para quem vai trabalhar na zona rural, a compra de casa simples no centro é preferível ao aluguel pelo baixo custo e praticidade.

\textbf{Áreas recomendadas:}
A área central da cidade de Caazapá é a única zona com oferta imobiliária minimamente estruturada no departamento, concentrando os poucos imóveis com escritura pública e acesso a serviços básicos.

\textbf{Armadilhas comuns:}
O mercado imobiliário é predominantemente informal, com transações feitas por contrato particular sem registro em cartório. Exija sempre escritura pública registrada e desconfie de preços muito abaixo da média regional, que podem indicar pendências jurídicas ou litigância sobre o imóvel.

\subsection*{Cobertura Celular}

{\small
{\small
\begin{tabularx}{\linewidth}{lXXXXX}
\toprule
Operadora & 2G & 3G & 4G & 5G & Cobertura Pop. \\
\midrule
Tigo & 78\% & 65\% & 52\% & --- & 70\% \\
Personal/Claro & 68\% & 55\% & 42\% & --- & 60\% \\
VOX & 40\% & 25\% & 12\% & --- & 35\% \\
\bottomrule
\end{tabularx}
}
}

> \textbf{Nota:} Caazapá é um dos departamentos com menor densidade populacional e menor cobertura telecomunicacional do Paraguai Oriental. A capital Caazapá e municípios da Ruta 8 têm cobertura razoável; o interior tem lacunas significativas.

\begin{tabularx}{\linewidth}{lX}
\toprule
Parâmetro & Valor \\
\midrule
Melhor operadora rural & Tigo \\
Velocidade 4G média (Mbps) & 12--25 \\
Áreas sem cobertura & Interior norte e leste, areas de mata, assentamentos indígenas \\
Sinal em propriedade rural & ruim a moderado \\
\bottomrule
\end{tabularx}

\begin{tabularx}{\linewidth}{lXXX}
\toprule
Plano & Operadora & USD/mês & Dados \\
\midrule
Pré-pago básico (7 dias) & Tigo & ~2,50 & 3 GB \\
Pré-pago intermediário (30 dias) & Personal & ~5,00 & 8 GB \\
Pós-pago entrada & Tigo & ~12,00 & 15 GB + chamadas \\
Pós-pago intermediário & Personal/Claro & ~18,00 & 30 GB + chamadas \\
\bottomrule
\end{tabularx}

\subsection*{Internet}

\begin{tabularx}{\linewidth}{lX}
\toprule
Indicador & Valor \\
\midrule
\% população com internet & ~52\% \\
\% com banda larga fixa & ~10\% dos domicílios \\
Velocidade média download (Mbps) & 10--25 (urbano) / 2--8 (rural) \\
Velocidade média upload (Mbps) & 3--10 (urbano) \\
\bottomrule
\end{tabularx}

> Caazapá apresenta um dos menores índices de acesso à internet fixa do Paraguai Oriental, reflexo da baixa densidade populacional e do nível de desenvolvimento econômico do departamento.

\begin{tabularx}{\linewidth}{lXX}
\toprule
Tecnologia & Disponibilidade & Velocidade Típica \\
\midrule
Fibra ótica (FTTH) & Apenas Caazapá cidade & 30--100 Mbps \\
Cabo HFC & Não disponível & --- \\
Rádio/Wireless & Municípios com +2k hab. & 3--20 Mbps \\
ADSL (Copaco) & Municípios com centrais telefônicas & 1--6 Mbps \\
Satélite (Starlink) & Todo o departamento & 50--200 Mbps \\
\bottomrule
\end{tabularx}

\begin{tabularx}{\linewidth}{lXXX}
\toprule
Provedor & Tecnologia & Área de cobertura & Preço ref. (USD/mês) \\
\midrule
Tigo & Fibra/Rádio & Caazapá cidade + municípios maiores & 18--40 \\
Claro/Personal & Rádio & Caazapá cidade & 15--35 \\
Copaco & ADSL & Interior com rede legada & 10--22 \\
ISPs locais & Rádio wireless & Municípios menores & 8--18 \\
Starlink & Satélite LEO & Todo o departamento & ~44 + equip. ~370 \\
\bottomrule
\end{tabularx}

A cobertura de internet rural em Caazapá é \textbf{muito limitada} --- um dos piores quadros no Paraguai Oriental:

\textbf{Recomendação para imigrantes com propriedades em Caazapá:} Orçar Starlink como item obrigatório na instalação. Equip. ~USD 370 (pagamento único) + ~USD 44/mês garante conectividade de 50--200 Mbps em qualquer ponto do departamento.

\clearpage
\newpage
\secaoDiagnostico{de Mayo}{\mapaDistrito{0725}}{3 de Mayo é um centro de estabilidade agrícola no sul de Caazapá.
Oferece uma combinação sólida de segurança institucional, abundância de
recursos hídricos (Rio Pirapó) e solos vocacionados para a fruticultura.
Sua força reside na baixa visibilidade estratégica e na coesão social
comunitária, sendo um local ideal para quem busca um refúgio produtivo e
discreto, com autonomia alimentar e hídrica garantida.

\subsubsection{Recomendação
Técnica}

Altamente indicado para estabelecimentos de autossuficiência de
pequeno/\allowbreak médio porte com foco em fruticultura ou pecuária sustentável.
Requer atenção à cota de inundação das propriedades e autonomia em
saúde. O GSS de 7.3 reflete a força da segurança social e dos recursos
naturais da localidade. Referência atual de score: GSS 7.3.

}
\begin{table}[ht!]
\small \begin{tabular}{p{3.0cm}p{0.8cm}p{6.0cm}} \toprule
\textbf{Critério} & \textbf{Nota} & \textbf{Justificativa} \\ \midrule
A - Ameaças Estratégicas & 8.1 & Posicionamento estratégico discreto mas bem conectado; fora de eixos de conflito \\
B - Risco Social & 7.5 & População pequena e densidade controlada; baixo risco de agitação social \\
C - Riscos Naturais & 5.9 & Geologia estável; nota penalizada pelo risco moderado de inundação ribeirinha \\
D - Autossuficiência & 7.0 & Excelentes recursos: água abundante, solo fértil para citros e proteína animal \\
E - Institucional & 7.4 & Forte estabilidade institucional regional e baixa criminalidade histórica \\
\midrule
 \textbf{GSS FINAL} & \textbf{7.3} & \textbf{Seguro (Localidade recomendada para quem busca qualidade de vida rural com foco em fruticultura e abundância hídrica).} \\ \bottomrule
 \end{tabular}
\end{table}
\subsubsection{Vulnerabilidades e
Mitigação}\label{vuln-3_de_Mayo-vulnerabilidades-e-mitigauxe7uxe3o}

\begin{itemize}
\tightlist
\item
  \textbf{Isolamento de Saúde:} Distância de centros hospitalares
  cirúrgicos (Mitigação: manutenção de estoque médico básico e kit
  trauma local).
\item
  \textbf{Risco Hídrico:} Inundações do Rio Pirapó (Mitigação:
  construção em cotas de segurança elevadas e preservação de matas
  ciliares).
\item
  \textbf{Instabilidade Elétrica:} Dependência da rede ANDE\footnote{\fineANDE} rural
  (Mitigação: instalação de sistemas solares fotovoltaicos redundantes).
\end{itemize}

\subsubsection{Dossiê de Campo}\label{dossie-3_de_Mayo-dossiuxea-de-campo}
\begin{description}[style=nextline,leftmargin=0.5cm,font=\bfseries]
\item[Coordenadas]  26°43'S 55°41'W.
\item[Topografia]  Relevo predominantemente plano com algumas ondulações. Localizado no sul do departamento, banhado pelo Rio Pirapó. Altitude \textasciitilde 140m. Geologicamente muito estável, risco sísmico desprezível.
\item[Alvos Estratégicos]  Áreas de produção citrícola (Laranja/\allowbreak Tangerina); Ponto de conexão rodoviária entre Caazapá e Itapúa. Ausência de alvos militares primários, oferecendo boa invisibilidade tática.
\item[Fallout]  Ventos predominantes de Nordeste (NE) e Leste (E). Baixo risco de fallout direto.
\item[População]  20.514 habitantes (Censo 2022 Final). Perfil majoritariamente rural com forte cultura de agricultura familiar.
\item[Densidade]  \textasciitilde 27,6 hab/\allowbreak km² (Baixa densidade, oferecendo excelente isolamento relativo em áreas de produção agrícola).
\item[Vias de Acesso]  Acesso via ramais asfálticos e estradas vicinais de terra batida. Obras recentes de pavimentação no trecho Lima - Tatú Kua melhoraram a logística regional. Risco de isolamento parcial em eventos de chuva extrema (>1.600 mm/\allowbreak ano).
\item[Serviços]  Unidade de Saúde da Família (USF) local. Dependência de Yuty ou da capital Caazapá para atendimentos médicos de alta complexidade.
\item[Custo de Vida]  Muito baixo; economia baseada na produção primária e pequena agroindústria.
\item[Preço da Terra]  US\$ 2.500-4.500/\allowbreak ha (áreas mistas/\allowbreak pecuária); US\$ 5.000-7.500/\allowbreak ha (áreas mecanizáveis ou com citricultura).
\item[Hidrologia]  Risco Moderado de Inundação. Proximidade com o Rio Pirapó e diversos arroios. Áreas ribeirinhas baixas vulneráveis a transbordamentos cíclicos em anos de El Niño intenso.
\item[Clima]  Subtropical Úmido. Verões quentes; invernos com frentes frias do Sul.
\item[Energia]  Rede nacional (Itaipu); instável em áreas rurais isoladas.
\item[Água]  Abundante via Rio Pirapó, arroios e poços artesianos; lençol freático produtivo.
\item[Qualidade do Solo]  Latossolo Vermelho e solos hidromórficos (nas baixadas); muito fértil em áreas altas, excelente para citros, mandioca e pastagens.
\item[Recursos Locais]  Grande produtor de laranjas e tangerinas, mandioca e gado de corte. Elevado potencial para autossuficiência alimentar básica em nível familiar e excedente de fruticultura.
\item[Segurança]  Zona rural pacífica. Caazapá mantém índices de segurança estáveis, embora enfrente desafios de microtráfico em núcleos urbanos maiores. Coesão social baseada em comunidades agrícolas tradicionais.
\item[Leis Local: Município consolidado. Livre de Restrição de Fronteira]  Localizado no interior do país, permitindo plena titularidade para brasileiros.
\end{description}
\textbf{Fonte climática:} NASA POWER Climatology API (período 2001-2020)

\textbf{Fonte luminosa:} estimativa world atlas

\paragraph{Irradiação Solar
(kWh/\allowbreak m²/\allowbreak dia)}\label{irradiauxe7uxe3o-solar-kwhmuxb2dia}


\begin{center}
\resizebox{\linewidth}{!}{%
{\renewcommand{\arraystretch}{1.2}%
\begin{tabular}{@{}c|c|c|c|c|c|c|c|c|c|c|c|c@{}}
\toprule
Jan & Fev & Mar & Abr & Mai & Jun & Jul & Ago & Set & Out & Nov & Dez & Média \\
\midrule
6.71 & 6.20 & 5.39 & 4.34 & 3.26 & 2.77 & 3.10 & 3.79 & 4.42 & 5.34 &
6.39 & 6.76 & \textbf{4.87} \\
\bottomrule
\end{tabular}}%
}
\end{center}


\textbf{Inclinação solar recomendada:} 26° N (anual) · 36° N (inverno
jun-ago) · 16° N (verão nov-jan)

\paragraph{Precipitação (mm/\allowbreak mês)}\label{precipitauxe7uxe3o-mmmuxeas}


\begin{center}
\resizebox{\linewidth}{!}{%
{\renewcommand{\arraystretch}{1.2}%
\begin{tabular}{@{}c|c|c|c|c|c|c|c|c|c|c|c|c@{}}
\toprule
Jan & Fev & Mar & Abr & Mai & Jun & Jul & Ago & Set & Out & Nov & Dez & Total/\allowbreak ano \\
\midrule
141 & 123 & 136 & 171 & 156 & 100 & 82 & 71 & 100 & 194 & 188 & 167 &
\textbf{1629 mm} \\
\bottomrule
\end{tabular}}%
}
\end{center}


\paragraph{Poluição Luminosa}\label{poluiuxe7uxe3o-luminosa}

\begin{longtable}[]{@{}ll@{}}
\toprule\noalign{}
Parâmetro & Valor \\
\midrule\noalign{}
\endhead
\bottomrule\noalign{}
\endlastfoot
Escala Bortle & 3 --- Céu rural \\
Radiância artificial & 0.8 nW/\allowbreak cm²/\allowbreak sr \\
\end{longtable}
\paragraph{Solo (SoilGrids 2.0, média ponderada 0--30
cm)}

\begin{longtable}[]{@{}ll@{}}
\toprule\noalign{}
Parâmetro & Valor \\
\midrule\noalign{}
\endhead
\bottomrule\noalign{}
\endlastfoot
pH (H$_2$O) & 5.2 \\
Carbono orgânico (SOC) & 21.32 g/\allowbreak kg \\
Argila & 35.02 \% \\
Areia & 30.12 \% \\
Silte (calc.) & 34.9 \% \\
Densidade aparente & 1.28 g/\allowbreak cm³ \\
\textbf{Aptidão agrícola} & \textbf{Média} \\
\end{longtable}

\subsubsection{Indicadores Sociais}

\textbf{Fontes:} PNUD 2020, INE\footnote{\fineINE} Censo 2022, INE\footnote{\fineINE} EPHC 2023, Ministerio
Público 2024\\
\textbf{Nota:} Valores marcados como \emph{(dept.)} referem-se ao
departamento de Caazapá; valores \emph{(dist.)} são específicos deste
distrito. Dados marcados \emph{(est.)} são estimativas.

\begin{longtable}[]{@{}
  >{\raggedright\arraybackslash}p{(\linewidth - 4\tabcolsep) * \real{0.4231}}
  >{\raggedright\arraybackslash}p{(\linewidth - 4\tabcolsep) * \real{0.2692}}
  >{\raggedright\arraybackslash}p{(\linewidth - 4\tabcolsep) * \real{0.3077}}@{}}
\toprule\noalign{}
\begin{minipage}[b]{\linewidth}\raggedright
Indicador
\end{minipage} & \begin{minipage}[b]{\linewidth}\raggedright
Valor
\end{minipage} & \begin{minipage}[b]{\linewidth}\raggedright
Âmbito
\end{minipage} \\
\midrule\noalign{}
\endhead
\bottomrule\noalign{}
\endlastfoot
IDH (2020) & 0,648 & dept. \\
Ranking IDH nacional & 18/\allowbreak 18 & dept. \\
Esperança de vida & 69,5 anos & dept. \\
Escolaridade média & 7.0 anos & dist. \\
RNB per capita (USD PPA) & 5.200 & dept. \\
Pobreza monetária (\%) & 34,7\% & dept. \\
Pobreza extrema (\%) & 11,3\% & dept. \\
Índice de Gini & 0,472 & dept. \\
Acesso a água potável (\%) & 60,1\% & dept. \\
Acesso a saneamento (\%) & 57,4\% & dept. \\
Taxa de homicídios (est., /\allowbreak 100k hab) & \textasciitilde4--6 (estimativa,
próxima à média nacional) & dept. \\
Índice de segurança & médio & dept. \\
População (Censo 2022) & 12.923 hab. & dist. \\
Idade mediana & 28.0 anos & dist. \\
Taxa de fecundidade & N/\allowbreak D & dist. \\
\end{longtable}
\subsubsection{Combustível}

\textbf{Referência:} PETROPAR /\allowbreak  postos locais (2024) \textbf{Tipo de
localidade:} Interior

\begin{longtable}[]{@{}lll@{}}
\toprule\noalign{}
Combustível & USD/\allowbreak litro & Gs/\allowbreak litro (aprox.) \\
\midrule\noalign{}
\endhead
\bottomrule\noalign{}
\endlastfoot
Gasolina 93 oct & 0.99 & 7,326 \\
Gasolina 97 oct (premium) & 1.09 & 8,066 \\
Diesel & 0.91 & 6,734 \\
\end{longtable}

\begin{quote}
Preços podem variar ±5\% conforme posto e sazonalidade. Chaco e interior
remoto apresentam maior variação.
\end{quote}

\subsubsection{Cobertura Celular}

\textbf{Fonte:} CONATEL PY /\allowbreak  operadoras (2024)

\begin{longtable}[]{@{}ll@{}}
\toprule\noalign{}
Parâmetro & Valor \\
\midrule\noalign{}
\endhead
\bottomrule\noalign{}
\endlastfoot
Cobertura 4G população (dept.) & 78\% \\
Cobertura 4G área rural & 55\% \\
Melhor operadora & Tigo \\
Qualidade rural & moderada \\
\end{longtable}

\begin{quote}
Para áreas rurais fora do núcleo urbano, recomenda-se chip Tigo como
principal e Personal como backup.
\end{quote}

\subsubsection{Internet}

\textbf{Fonte:} CONATEL /\allowbreak  Speedtest Ookla (2024)

\begin{longtable}[]{@{}ll@{}}
\toprule\noalign{}
Parâmetro & Valor \\
\midrule\noalign{}
\endhead
\bottomrule\noalign{}
\endlastfoot
Velocidade média download & 32 Mbps \\
Domicílios com internet (dept.) & 42\% \\
Tecnologia predominante & rádio \\
Opção rural & Starlink disponível (\textasciitilde USD 44/\allowbreak mês) \\
\end{longtable}

\subsubsection{Mercado Imobiliário e Terra
Rural}

\textbf{Fonte:} INDERT /\allowbreak  Clasificados.com.py (2024)

\begin{longtable}[]{@{}ll@{}}
\toprule\noalign{}
Tipo & Referência \\
\midrule\noalign{}
\endhead
\bottomrule\noalign{}
\endlastfoot
Terra agrícola alta prod. (USD/\allowbreak ha) & 5,500 \\
Imóvel urbano (USD/\allowbreak m²) & 550 \\
Aluguel 2 quartos (USD/\allowbreak mês) & 220 \\
\end{longtable}

\begin{quote}
Valores de referência departamental. Localidades menores podem ter
preços 20--40\% abaixo da capital departamental. \#\#\# Saúde
\end{quote}

\textbf{Fonte:} MSPBS /\allowbreak  IPS Paraguay (2024-2026), consolidação
departamental e proxy local

\begin{longtable}[]{@{}ll@{}}
\toprule\noalign{}
Serviço & Disponibilidade \\
\midrule\noalign{}
\endhead
\bottomrule\noalign{}
\endlastfoot
USF /\allowbreak  Posto de Saúde & sim \\
Hospital Regional & não \\
IPS (seguro social) & não \\
Farmácia & sim \\
Distância ao hospital de referência & \textasciitilde52 km (Caazapa) \\
\end{longtable}

\textbf{Principais estabelecimentos:} USF local; referencia hospitalar
em Caazapa

\textbf{Observação para imigrantes:} Atendimento primario local ou em
raio curto; casos de maior complexidade seguem para Caazapa. Cobertura
privada continua recomendada para especialidades e urgências de maior
complexidade.
\newpage
\secaoDiagnostico{Abai}{\mapaDistrito{0602}}{Abaí é a nova fronteira de prosperidade do departamento de Caazapá.
Oferece uma combinação sólida de segurança institucional, abundância de
recursos agropecuários e solos de elite (``terra vermelha''). Sua
localização estratégica na nova rota de integração para o Leste garante
excelente conectividade comercial sem comprometer o isolamento tático
provido pela sua vasta extensão territorial. É o local ideal para quem
busca produtividade agrícola de larga escala e segurança social.

\subsubsection{Recomendação
Técnica}

Altamente indicado para estabelecimentos de autossuficiência comercial
ou agroindústria de grãos. Requer vigilância jurídica sobre a
titularidade das terras e autonomia em saúde básica. O GSS de 7.4
reflete a força econômica e a segurança da localidade. Referência atual
de score: GSS 7.4.

}
\begin{table}[ht!]
\small \begin{tabular}{p{3.0cm}p{0.8cm}p{6.0cm}} \toprule
\textbf{Critério} & \textbf{Nota} & \textbf{Justificativa} \\ \midrule
A - Ameaças Estratégicas & 8.2 & Posicionamento estratégico tático como hub de integração para o Leste; fora de eixos de conflito primário \\
B - Risco Social & 7.5 & População produtiva dispersa em território vasto; isolamento demográfico favorável \\
C - Riscos Naturais & 5.8 & Geologia muito estável; nota penalizada por riscos hidrometeorológicos moderados e erosão \\
D - Autossuficiência & 7.5 & Recursos agropecuários excepcionais: solo de elite, grãos e água abundante \\
E - Institucional & 7.4 & Forte estabilidade institucional regional e baixa criminalidade histórica \\
\midrule
 \textbf{GSS FINAL} & \textbf{7.4} & \textbf{Seguro (Localidade altamente recomendada para quem busca alta produtividade agrícola em ambiente seguro).} \\ \bottomrule
 \end{tabular}
\end{table}
\subsubsection{Vulnerabilidades e
Mitigação}\label{vuln-Abai-vulnerabilidades-e-mitigauxe7uxe3o}

\begin{itemize}
\tightlist
\item
  \textbf{Conflitos de Terra:} Risco de invasões em grandes propriedades
  (Mitigação: aquisição de terras com título perfeito, integração social
  local e segurança privada patrimonial).
\item
  \textbf{Isolamento de Saúde:} Distância de centros hospitalares
  cirúrgicos (Mitigação: manutenção de estoque médico básico e
  treinamento de socorrista).
\item
  \textbf{Dependência Logística:} Dependência da nova rota asfáltica
  para escoamento (Mitigação: estoque de insumos e combustíveis para 60
  dias).
\end{itemize}

\subsubsection{Dossiê de Campo}\label{dossie-Abai-dossiuxea-de-campo}
\begin{description}[style=nextline,leftmargin=0.5cm,font=\bfseries]
\item[Coordenadas]  26°02'S 55°56'W.
\item[Topografia]  Relevo ondulado com extensas planícies mecanizáveis e áreas de mata nativa preservada. Altitude \textasciitilde 200m. Área vasta de \textasciitilde 1.922 km². Geologicamente muito estável, risco sísmico desprezível.
\item[Alvos Estratégicos]  Áreas de expansão acelerada de soja e milho (um dos maiores produtores departamentais); Ponto de conexão da "Rota da Integração" ligando Caazapá ao Alto Paraná (San Cristóbal). Ausência de alvos militares primários, oferecendo boa invisibilidade estratégica em território vasto.
\item[Fallout]  Ventos predominantes do Norte (quentes e úmidos) e Leste. Baixo risco de fallout direto.
\item[População]  38.589 habitantes (Censo 2022 Final). Distrito com forte componente demográfico de colonos brasileiros e paraguaios focados no agronegócio.
\item[Densidade]  \textasciitilde 20,1 hab/\allowbreak km² (Baixa densidade, considerando a grande extensão territorial, oferecendo isolamento tático e privacidade estratégica).
\item[Vias de Acesso]  Acesso via ramais asfálticos recentes (Rota da Integração) que melhoraram drasticamente a conexão com Ciudad del Este e polos de Alto Paraná. Conectividade interna por estradas de terra que exigem veículos robustos em períodos chuvosos (média 1.700 mm/\allowbreak ano).
\item[Serviços]  Centro de Saúde local e serviços de apoio às cooperativas agrícolas. Dependência de San Juan Nepomuceno ou Villarrica para saúde de alta complexidade.
\item[Custo de Vida]  Baixo a Médio; economia dolarizada pelo agronegócio exportador.
\item[Preço da Terra]  US\$ 5.000-8.500/\allowbreak ha (terras mecanizadas de elite); US\$ 2.500-4.000/\allowbreak ha (áreas de campo natural/\allowbreak pecuária).
\item[Hidrologia]  Diversos arroios locais integrados a sistemas de drenagem agrícola. Baixo risco de inundações sistêmicas catastróficas. Risco moderado de erosão em solos descobertos durante tempestades severas.
\item[Clima]  Subtropical Úmido. Verões intensos e invernos com geadas ocasionais.
\item[Energia]  Rede nacional (Itaipu); instável em áreas rurais remotas durante picos de verão.
\item[Água]  Abundante via poços artesianos e lençol freático produtivo (influência da bacia do Rio Monday ao norte).
\item[Qualidade do Solo]  Latossolo Vermelho (Terra Vermelha); profundo, fértil e mecanizável. Excelente aptidão para soja, milho, trigo e gado de engorda.
\item[Recursos Locais]  Grande produção de grãos e proteína animal. Elevadíssimo potencial para autossuficiência alimentar e hídrica em nível de propriedade.
\item[Segurança]  Zona rural tradicionalmente pacífica, mas sob vigilância tática devido à riqueza do agronegócio. Baixa criminalidade urbana. Coesão social baseada na cultura produtora. Risco pontual de conflitos agrários e invasões em áreas de mata.
\item[Leis Local: Município consolidado e polo de influência econômica. Livre de Restrição de Fronteira]  Fora da faixa de 50km da fronteira, permitindo plena titularidade para brasileiros.
\end{description}
\textbf{Fonte climática:} NASA POWER Climatology API (período 2001-2020)

\textbf{Fonte luminosa:} estimativa world atlas

\paragraph{Irradiação Solar
(kWh/\allowbreak m²/\allowbreak dia)}\label{irradiauxe7uxe3o-solar-kwhmuxb2dia}


\begin{center}
\resizebox{\linewidth}{!}{%
{\renewcommand{\arraystretch}{1.2}%
\begin{tabular}{@{}c|c|c|c|c|c|c|c|c|c|c|c|c@{}}
\toprule
Jan & Fev & Mar & Abr & Mai & Jun & Jul & Ago & Set & Out & Nov & Dez & Média \\
\midrule
6.56 & 6.06 & 5.32 & 4.33 & 3.21 & 2.76 & 3.08 & 3.81 & 4.39 & 5.25 &
6.31 & 6.67 & \textbf{4.81} \\
\bottomrule
\end{tabular}}%
}
\end{center}


\textbf{Inclinação solar recomendada:} 26° N (anual) · 36° N (inverno
jun-ago) · 16° N (verão nov-jan)

\paragraph{Precipitação (mm/\allowbreak mês)}\label{precipitauxe7uxe3o-mmmuxeas}


\begin{center}
\resizebox{\linewidth}{!}{%
{\renewcommand{\arraystretch}{1.2}%
\begin{tabular}{@{}c|c|c|c|c|c|c|c|c|c|c|c|c@{}}
\toprule
Jan & Fev & Mar & Abr & Mai & Jun & Jul & Ago & Set & Out & Nov & Dez & Total/\allowbreak ano \\
\midrule
135 & 129 & 133 & 159 & 166 & 94 & 78 & 61 & 93 & 189 & 188 & 164 &
\textbf{1590 mm} \\
\bottomrule
\end{tabular}}%
}
\end{center}


\paragraph{Poluição Luminosa}\label{poluiuxe7uxe3o-luminosa}

\begin{longtable}[]{@{}ll@{}}
\toprule\noalign{}
Parâmetro & Valor \\
\midrule\noalign{}
\endhead
\bottomrule\noalign{}
\endlastfoot
Escala Bortle & 3 --- Céu rural \\
Radiância artificial & 0.8 nW/\allowbreak cm²/\allowbreak sr \\
\end{longtable}
\paragraph{Solo (SoilGrids 2.0, média ponderada 0--30
cm)}

\begin{longtable}[]{@{}ll@{}}
\toprule\noalign{}
Parâmetro & Valor \\
\midrule\noalign{}
\endhead
\bottomrule\noalign{}
\endlastfoot
pH (H$_2$O) & 5.3 \\
Carbono orgânico (SOC) & 18.17 g/\allowbreak kg \\
Argila & 29.45 \% \\
Areia & 38.66 \% \\
Silte (calc.) & 31.9 \% \\
Densidade aparente & 1.32 g/\allowbreak cm³ \\
\textbf{Aptidão agrícola} & \textbf{Média} \\
\end{longtable}

\begin{itemize}
\tightlist
\item
  \textbf{Homicidios/\allowbreak 100k:} \textasciitilde7,0 (Regional/\allowbreak Caazapá).
\end{itemize}
\subsubsection{Indicadores Sociais}

\textbf{Fontes:} PNUD 2020, INE\footnote{\fineINE} Censo 2022, INE\footnote{\fineINE} EPHC 2023, Ministerio
Público 2024\\
\textbf{Nota:} Valores marcados como \emph{(dept.)} referem-se ao
departamento de Caazapá; valores \emph{(dist.)} são específicos deste
distrito. Dados marcados \emph{(est.)} são estimativas.

\begin{longtable}[]{@{}
  >{\raggedright\arraybackslash}p{(\linewidth - 4\tabcolsep) * \real{0.4231}}
  >{\raggedright\arraybackslash}p{(\linewidth - 4\tabcolsep) * \real{0.2692}}
  >{\raggedright\arraybackslash}p{(\linewidth - 4\tabcolsep) * \real{0.3077}}@{}}
\toprule\noalign{}
\begin{minipage}[b]{\linewidth}\raggedright
Indicador
\end{minipage} & \begin{minipage}[b]{\linewidth}\raggedright
Valor
\end{minipage} & \begin{minipage}[b]{\linewidth}\raggedright
Âmbito
\end{minipage} \\
\midrule\noalign{}
\endhead
\bottomrule\noalign{}
\endlastfoot
IDH (2020) & 0,648 & dept. \\
Ranking IDH nacional & 18/\allowbreak 18 & dept. \\
Esperança de vida & 69,5 anos & dept. \\
Escolaridade média & 7.3 anos & dist. \\
RNB per capita (USD PPA) & 5.200 & dept. \\
Pobreza monetária (\%) & 34,7\% & dept. \\
Pobreza extrema (\%) & 11,3\% & dept. \\
Índice de Gini & 0,472 & dept. \\
Acesso a água potável (\%) & 60,1\% & dept. \\
Acesso a saneamento (\%) & 57,4\% & dept. \\
Taxa de homicídios (est., /\allowbreak 100k hab) & \textasciitilde4--6 (estimativa,
próxima à média nacional) & dept. \\
Índice de segurança & médio & dept. \\
População (Censo 2022) & 21.964 hab. & dist. \\
Idade mediana & 26.0 anos & dist. \\
Taxa de fecundidade & N/\allowbreak D & dist. \\
\end{longtable}
\subsubsection{Combustível}

\textbf{Referência:} PETROPAR /\allowbreak  postos locais (2024) \textbf{Tipo de
localidade:} Interior

\begin{longtable}[]{@{}lll@{}}
\toprule\noalign{}
Combustível & USD/\allowbreak litro & Gs/\allowbreak litro (aprox.) \\
\midrule\noalign{}
\endhead
\bottomrule\noalign{}
\endlastfoot
Gasolina 93 oct & 0.99 & 7,326 \\
Gasolina 97 oct (premium) & 1.09 & 8,066 \\
Diesel & 0.91 & 6,734 \\
\end{longtable}

\begin{quote}
Preços podem variar ±5\% conforme posto e sazonalidade. Chaco e interior
remoto apresentam maior variação.
\end{quote}

\subsubsection{Cobertura Celular}

\textbf{Fonte:} CONATEL PY /\allowbreak  operadoras (2024)

\begin{longtable}[]{@{}ll@{}}
\toprule\noalign{}
Parâmetro & Valor \\
\midrule\noalign{}
\endhead
\bottomrule\noalign{}
\endlastfoot
Cobertura 4G população (dept.) & 78\% \\
Cobertura 4G área rural & 55\% \\
Melhor operadora & Tigo \\
Qualidade rural & moderada \\
\end{longtable}

\begin{quote}
Para áreas rurais fora do núcleo urbano, recomenda-se chip Tigo como
principal e Personal como backup.
\end{quote}

\subsubsection{Internet}

\textbf{Fonte:} CONATEL /\allowbreak  Speedtest Ookla (2024)

\begin{longtable}[]{@{}ll@{}}
\toprule\noalign{}
Parâmetro & Valor \\
\midrule\noalign{}
\endhead
\bottomrule\noalign{}
\endlastfoot
Velocidade média download & 32 Mbps \\
Domicílios com internet (dept.) & 42\% \\
Tecnologia predominante & rádio \\
Opção rural & Starlink disponível (\textasciitilde USD 44/\allowbreak mês) \\
\end{longtable}

\subsubsection{Mercado Imobiliário e Terra
Rural}

\textbf{Fonte:} INDERT /\allowbreak  Clasificados.com.py (2024)

\begin{longtable}[]{@{}ll@{}}
\toprule\noalign{}
Tipo & Referência \\
\midrule\noalign{}
\endhead
\bottomrule\noalign{}
\endlastfoot
Terra agrícola alta prod. (USD/\allowbreak ha) & 5,500 \\
Imóvel urbano (USD/\allowbreak m²) & 550 \\
Aluguel 2 quartos (USD/\allowbreak mês) & 220 \\
\end{longtable}

\begin{quote}
Valores de referência departamental. Localidades menores podem ter
preços 20--40\% abaixo da capital departamental. \#\#\# Saúde
\end{quote}

\textbf{Fonte:} MSPBS /\allowbreak  IPS Paraguay (2024-2026), consolidação
departamental e proxy local

\begin{longtable}[]{@{}ll@{}}
\toprule\noalign{}
Serviço & Disponibilidade \\
\midrule\noalign{}
\endhead
\bottomrule\noalign{}
\endlastfoot
USF /\allowbreak  Posto de Saúde & sim \\
Hospital Regional & não \\
IPS (seguro social) & não \\
Farmácia & sim \\
Distância ao hospital de referência & \textasciitilde57 km (Caazapa) \\
\end{longtable}

\textbf{Principais estabelecimentos:} USF local; referencia hospitalar
em Caazapa

\textbf{Observação para imigrantes:} Atendimento primario local ou em
raio curto; casos de maior complexidade seguem para Caazapa. Cobertura
privada continua recomendada para especialidades e urgências de maior
complexidade.
\newpage
\secaoDiagnostico{Buena Vista}{\mapaDistrito{0607}}{Buena Vista é um refúgio de produtividade e paz social no centro de
Caazapá. Sua maior força reside na combinação de baixa visibilidade
tática com uma segurança comunitária robusta e solos de alta qualidade
para a fruticultura. Oferece um santuário rural discreto, com fácil
acesso aos serviços da capital departamental, sendo ideal para projetos
de autossuficiência familiar que priorizam a estabilidade e a qualidade
de vida.

\subsubsection{Recomendação
Técnica}

Altamente indicado para estabelecimentos de autossuficiência de
pequeno/\allowbreak médio porte que valorizem a privacidade demográfica e a
segurança institucional superior da região. Requer autonomia em
logística básica e saúde. O GSS de 6.7 reflete a solidez da base de
recursos e a segurança social da localidade. Referência atual de score:
GSS 6.7.

}
\begin{table}[ht!]
\small \begin{tabular}{p{3.0cm}p{0.8cm}p{6.0cm}} \toprule
\textbf{Critério} & \textbf{Nota} & \textbf{Justificativa} \\ \midrule
A - Ameaças Estratégicas & 6.5 & Localização discreta e protegida de rotas principais; excelente invisibilidade tática \\
B - Risco Social & 7.0 & População reduzida e densidade controlada; baixíssimo risco social urbano \\
C - Riscos Naturais & 7.0 & Geologia muito estável e riscos naturais climáticos típicos manejáveis \\
D - Autossuficiência & 7.0 & Excelentes recursos: solo fértil para fruticultura, abundância hídrica e proteína animal \\
E - Institucional & 6.0 & Ambiente institucional seguro e pacífico, mas com infraestrutura pública local limitada \\
\midrule
 \textbf{GSS FINAL} & \textbf{6.7} & \textbf{Moderadamente Seguro (Ideal para quem busca refúgio rural tranquilo com facilidade de produção alimentar).} \\ \bottomrule
 \end{tabular}
\end{table}
\subsubsection{Vulnerabilidades e
Mitigação}\label{vuln-Buena_Vista-vulnerabilidades-e-mitigauxe7uxe3o}

\begin{itemize}
\tightlist
\item
  \textbf{Isolamento de Saúde:} Necessidade de deslocamento para Caazapá
  ou Villarrica em emergências (Mitigação: manutenção de estoque médico
  básico e kit trauma local).
\item
  \textbf{Dependência Econômica Citrícola:} Sensibilidade do setor a
  pragas ou crises de mercado (Mitigação: diversificação com horta de
  subsistência e criação de pequenos animais).
\item
  \textbf{Instabilidade Elétrica Rural:} Quedas ocasionais em temporais
  (Mitigação: instalação de sistemas solares fotovoltaicos redundantes).
\end{itemize}

\subsubsection{Dossiê de Campo}\label{dossie-Buena_Vista-dossiuxea-de-campo}
\begin{description}[style=nextline,leftmargin=0.5cm,font=\bfseries]
\item[Coordenadas]  26°16'S 56°01'W.
\item[Topografia]  Relevo predominantemente plano a suavemente ondulado. Localizado na zona central do departamento. Altitude \textasciitilde 160m. Área de \textasciitilde 126 km². Geologicamente muito estável, risco sísmico desprezível.
\item[Alvos Estratégicos]  Áreas de produção citrícola e pecuária; Ponto de passagem secundário ligando a capital Caazapá a San Juan Nepomuceno. Ausência de alvos militares ou industriais primários, oferecendo excelente invisibilidade tática e discrição estratégica.
\item[Fallout]  Ventos predominantes do Norte (NE) e Leste. Baixo risco de fallout direto.
\item[População]  6.319 habitantes (Censo 2022 Final). Perfil demográfico equilibrado entre núcleos povoados e áreas rurais.
\item[Densidade]  \textasciitilde 50,1 hab/\allowbreak km² (Densidade moderada-baixa, oferecendo isolamento relativo em propriedades rurais integradas).
\item[Vias de Acesso]  Acesso via ramal asfáltico que conecta à Ruta PY08. Localização estratégica "fora do eixo" de grandes rodovias nacionais, servindo como um porto seguro rural.
\item[Serviços]  Unidade de Saúde da Família (USF) local. Dependência direta da capital Caazapá (15 km) para serviços médicos de alta complexidade e logística comercial diversificada.
\item[Custo de Vida]  Muito baixo; economia baseada na citricultura familiar e pecuária de pequena/\allowbreak média escala.
\item[Preço da Terra]  US\$ 3.000-5.500/\allowbreak ha (áreas rurais agrícolas); US\$ 1.500-2.500/\allowbreak ha (áreas de campo natural/\allowbreak mata).
\item[Hidrologia]  Baixo risco de inundações sistêmicas catastróficas. Presença de arroios perenes favoráveis à pequena agricultura e pecuária. Drenagem natural eficiente.
\item[Clima]  Subtropical Úmido. Verões quentes; média pluvial de 1.600 mm/\allowbreak ano bem distribuída.
\item[Energia]  Rede nacional (Itaipu); boa estabilidade devido à proximidade com o hub regional da capital.
\item[Água]  Abundante via poços artesianos e juntas de saneamento; boa disponibilidade hídrica subterrânea.
\item[Qualidade do Solo]  Latossolo Vermelho e solos podzólicos; férteis e profundos, excelentes para citros (laranja/\allowbreak tangerina), mandioca e pastagens.
\item[Recursos Locais]  Grande produção de citros, mel, mandioca e gado de corte. Elevado potencial para autossuficiência alimentar básica em nível familiar.
\item[Segurança]  Zona extremamente pacífica e estável. Taxa de homicídios no departamento de Caazapá entre as menores do país (\textasciitilde 7,0/\allowbreak 100k). Coesão social forte baseada em tradições rurais conservadoras. Baixíssima incidência de crimes violentos graves.
\item[Leis Local: Município consolidado. Livre de Restrição de Fronteira]  Localizado no interior do país, permitindo plena titularidade para investidores brasileiros.
\end{description}
\textbf{Fonte climática:} NASA POWER Climatology API (período 2001-2020)

\textbf{Fonte luminosa:} estimativa world atlas

\paragraph{Irradiação Solar
(kWh/\allowbreak m²/\allowbreak dia)}\label{irradiauxe7uxe3o-solar-kwhmuxb2dia}


\begin{center}
\resizebox{\linewidth}{!}{%
{\renewcommand{\arraystretch}{1.2}%
\begin{tabular}{@{}c|c|c|c|c|c|c|c|c|c|c|c|c@{}}
\toprule
Jan & Fev & Mar & Abr & Mai & Jun & Jul & Ago & Set & Out & Nov & Dez & Média \\
\midrule
6.71 & 6.20 & 5.39 & 4.34 & 3.26 & 2.77 & 3.10 & 3.79 & 4.42 & 5.34 &
6.39 & 6.76 & \textbf{4.87} \\
\bottomrule
\end{tabular}}%
}
\end{center}


\textbf{Inclinação solar recomendada:} 26° N (anual) · 36° N (inverno
jun-ago) · 16° N (verão nov-jan)

\paragraph{Precipitação (mm/\allowbreak mês)}\label{precipitauxe7uxe3o-mmmuxeas}


\begin{center}
\resizebox{\linewidth}{!}{%
{\renewcommand{\arraystretch}{1.2}%
\begin{tabular}{@{}c|c|c|c|c|c|c|c|c|c|c|c|c@{}}
\toprule
Jan & Fev & Mar & Abr & Mai & Jun & Jul & Ago & Set & Out & Nov & Dez & Total/\allowbreak ano \\
\midrule
135 & 129 & 133 & 159 & 166 & 94 & 78 & 61 & 93 & 189 & 188 & 164 &
\textbf{1590 mm} \\
\bottomrule
\end{tabular}}%
}
\end{center}


\paragraph{Poluição Luminosa}\label{poluiuxe7uxe3o-luminosa}

\begin{longtable}[]{@{}ll@{}}
\toprule\noalign{}
Parâmetro & Valor \\
\midrule\noalign{}
\endhead
\bottomrule\noalign{}
\endlastfoot
Escala Bortle & 3 --- Céu rural \\
Radiância artificial & 0.8 nW/\allowbreak cm²/\allowbreak sr \\
\end{longtable}
\paragraph{Solo (SoilGrids 2.0, média ponderada 0--30
cm)}

\begin{longtable}[]{@{}ll@{}}
\toprule\noalign{}
Parâmetro & Valor \\
\midrule\noalign{}
\endhead
\bottomrule\noalign{}
\endlastfoot
pH (H$_2$O) & 5.3 \\
Carbono orgânico (SOC) & 18.3 g/\allowbreak kg \\
Argila & 28.9 \% \\
Areia & 39.92 \% \\
Silte (calc.) & 31.2 \% \\
Densidade aparente & 1.36 g/\allowbreak cm³ \\
\textbf{Aptidão agrícola} & \textbf{Média} \\
\end{longtable}

\subsubsection{Indicadores Sociais}

\textbf{Fontes:} PNUD 2020, INE\footnote{\fineINE} Censo 2022, INE\footnote{\fineINE} EPHC 2023, Ministerio
Público 2024\\
\textbf{Nota:} Valores marcados como \emph{(dept.)} referem-se ao
departamento de Caazapá; valores \emph{(dist.)} são específicos deste
distrito. Dados marcados \emph{(est.)} são estimativas.

\begin{longtable}[]{@{}
  >{\raggedright\arraybackslash}p{(\linewidth - 4\tabcolsep) * \real{0.4231}}
  >{\raggedright\arraybackslash}p{(\linewidth - 4\tabcolsep) * \real{0.2692}}
  >{\raggedright\arraybackslash}p{(\linewidth - 4\tabcolsep) * \real{0.3077}}@{}}
\toprule\noalign{}
\begin{minipage}[b]{\linewidth}\raggedright
Indicador
\end{minipage} & \begin{minipage}[b]{\linewidth}\raggedright
Valor
\end{minipage} & \begin{minipage}[b]{\linewidth}\raggedright
Âmbito
\end{minipage} \\
\midrule\noalign{}
\endhead
\bottomrule\noalign{}
\endlastfoot
IDH (2020) & 0,648 & dept. \\
Ranking IDH nacional & 18/\allowbreak 18 & dept. \\
Esperança de vida & 69,5 anos & dept. \\
Escolaridade média & 8.0 anos & dist. \\
RNB per capita (USD PPA) & 5.200 & dept. \\
Pobreza monetária (\%) & 34,7\% & dept. \\
Pobreza extrema (\%) & 11,3\% & dept. \\
Índice de Gini & 0,472 & dept. \\
Acesso a água potável (\%) & 60,1\% & dept. \\
Acesso a saneamento (\%) & 57,4\% & dept. \\
Taxa de homicídios (est., /\allowbreak 100k hab) & \textasciitilde4--6 (estimativa,
próxima à média nacional) & dept. \\
Índice de segurança & médio & dept. \\
População (Censo 2022) & 4.296 hab. & dist. \\
Idade mediana & 32.0 anos & dist. \\
Taxa de fecundidade & N/\allowbreak D & dist. \\
\end{longtable}
\subsubsection{Combustível}

\textbf{Referência:} PETROPAR /\allowbreak  postos locais (2024) \textbf{Tipo de
localidade:} Interior

\begin{longtable}[]{@{}lll@{}}
\toprule\noalign{}
Combustível & USD/\allowbreak litro & Gs/\allowbreak litro (aprox.) \\
\midrule\noalign{}
\endhead
\bottomrule\noalign{}
\endlastfoot
Gasolina 93 oct & 0.99 & 7,326 \\
Gasolina 97 oct (premium) & 1.09 & 8,066 \\
Diesel & 0.91 & 6,734 \\
\end{longtable}

\begin{quote}
Preços podem variar ±5\% conforme posto e sazonalidade. Chaco e interior
remoto apresentam maior variação.
\end{quote}

\subsubsection{Cobertura Celular}

\textbf{Fonte:} CONATEL PY /\allowbreak  operadoras (2024)

\begin{longtable}[]{@{}ll@{}}
\toprule\noalign{}
Parâmetro & Valor \\
\midrule\noalign{}
\endhead
\bottomrule\noalign{}
\endlastfoot
Cobertura 4G população (dept.) & 78\% \\
Cobertura 4G área rural & 55\% \\
Melhor operadora & Tigo \\
Qualidade rural & moderada \\
\end{longtable}

\begin{quote}
Para áreas rurais fora do núcleo urbano, recomenda-se chip Tigo como
principal e Personal como backup.
\end{quote}

\subsubsection{Internet}

\textbf{Fonte:} CONATEL /\allowbreak  Speedtest Ookla (2024)

\begin{longtable}[]{@{}ll@{}}
\toprule\noalign{}
Parâmetro & Valor \\
\midrule\noalign{}
\endhead
\bottomrule\noalign{}
\endlastfoot
Velocidade média download & 32 Mbps \\
Domicílios com internet (dept.) & 42\% \\
Tecnologia predominante & rádio \\
Opção rural & Starlink disponível (\textasciitilde USD 44/\allowbreak mês) \\
\end{longtable}

\subsubsection{Mercado Imobiliário e Terra
Rural}

\textbf{Fonte:} INDERT /\allowbreak  Clasificados.com.py (2024)

\begin{longtable}[]{@{}ll@{}}
\toprule\noalign{}
Tipo & Referência \\
\midrule\noalign{}
\endhead
\bottomrule\noalign{}
\endlastfoot
Terra agrícola alta prod. (USD/\allowbreak ha) & 5,500 \\
Imóvel urbano (USD/\allowbreak m²) & 550 \\
Aluguel 2 quartos (USD/\allowbreak mês) & 220 \\
\end{longtable}

\begin{quote}
Valores de referência departamental. Localidades menores podem ter
preços 20--40\% abaixo da capital departamental. \#\#\# Saúde
\end{quote}

\textbf{Fonte:} MSPBS /\allowbreak  IPS Paraguay (2024-2026), consolidação
departamental e proxy local

\begin{longtable}[]{@{}ll@{}}
\toprule\noalign{}
Serviço & Disponibilidade \\
\midrule\noalign{}
\endhead
\bottomrule\noalign{}
\endlastfoot
USF /\allowbreak  Posto de Saúde & sim \\
Hospital Regional & não \\
IPS (seguro social) & não \\
Farmácia & sim \\
Distância ao hospital de referência & \textasciitilde35 km (Caazapa) \\
\end{longtable}

\textbf{Principais estabelecimentos:} USF local; referencia hospitalar
em Caazapa

\textbf{Observação para imigrantes:} Atendimento primario local ou em
raio curto; casos de maior complexidade seguem para Caazapa. Cobertura
privada continua recomendada para especialidades e urgências de maior
complexidade.
\newpage
\secaoDiagnostico{Caazapa}{\mapaDistrito{0601}}{Caazapá é o porto seguro do interior paraguaio. Oferece a melhor
infraestrutura de serviços e segurança institucional do departamento,
garantindo um ambiente social de elite e baixa densidade populacional.
Sua força reside na abundância de recursos hídricos históricos e na
solidez de sua base institucional, sendo o local ideal para quem busca
proximidade com serviços de saúde e administrativos sem a densidade de
capitais maiores.

\subsubsection{Recomendação
Técnica}

Altamente indicado como base logística e administrativa para o projeto.
Requer posicionamento em zonas rurais periféricas para maximizar o
isolamento social. O GSS de 7.0 reflete a excepcional segurança física e
institucional da localidade. Referência atual de score: GSS 7.0.

}
\begin{table}[ht!]
\small \begin{tabular}{p{3.0cm}p{0.8cm}p{6.0cm}} \toprule
\textbf{Critério} & \textbf{Nota} & \textbf{Justificativa} \\ \midrule
A - Ameaças Estratégicas & 7.1 & Capital departamental e hub logístico; alvo estratégico tático administrativo seguro \\
B - Risco Social & 7.8 & Baixa densidade demográfica regional e população pequena; excelente segurança social \\
C - Riscos Naturais & 5.8 & Geologia estável; nota penalizada por riscos hidrometeorológicos típicos sazonais \\
D - Autossuficiência & 6.4 & Excelentes recursos hídricos e proteína animal; solo de elite mas limitado pela escala urbana \\
E - Institucional & 7.7 & Forte segurança institucional e serviços de saúde de referência regional \\
\midrule
 \textbf{GSS FINAL} & \textbf{7.0} & \textbf{Seguro (Ideal como base administrativa e residencial segura para o projeto).} \\ \bottomrule
 \end{tabular}
\end{table}
\subsubsection{Vulnerabilidades e
Mitigação}\label{vuln-Caazapa-vulnerabilidades-e-mitigauxe7uxe3o}

\begin{itemize}
\tightlist
\item
  \textbf{Pressão Administrativa:} Alvo burocrático em cenários de
  instabilidade política nacional (Mitigação: residência rural afastada
  5-10 km do centro administrativo).
\item
  \textbf{Dependência de Serviços Externos:} Para insumos industriais
  especializados (Mitigação: estoque de suprimentos técnicos para 60
  dias).
\item
  \textbf{Vulnerabilidade Climática:} Verões com temperaturas extremas
  (Mitigação: construção com alta inércia térmica e sistemas de
  resfriamento solar).
\end{itemize}

\subsubsection{Dossiê de Campo}\label{dossie-Caazapa-dossiuxea-de-campo}
\begin{description}[style=nextline,leftmargin=0.5cm,font=\bfseries]
\item[Coordenadas]  26°11'S 56°22'W.
\item[Topografia]  Relevo plano com suaves ondulações e vales férteis. Altitude \textasciitilde 155m. Geologicamente muito estável, risco sísmico praticamente nulo.
\item[Alvos Estratégicos]  Sede do Governo Departamental; Ponto de conexão da Ruta PY08 (Eixo Norte-Sul); Centro administrativo e de comando regional. Ausência de alvos militares de grande escala, oferecendo boa segurança institucional.
\item[Fallout]  Ventos predominantes do Norte (quentes) e Leste. Baixo risco de fallout direto.
\item[População]  14.322 habitantes (Censo 2022 Final). Perfil demográfico equilibrado com forte identidade cultural e histórica.
\item[Densidade]  \textasciitilde 17,4 hab/\allowbreak km² (Baixa densidade, oferecendo isolamento relativo com excelente suporte estatal local).
\item[Vias de Acesso]  Localização estratégica sobre a Ruta PY08 (asfaltada). Conectividade eficiente para o Norte (Villarrica) e Sul (Itapúa). Risco de saturação rodoviária moderado em crises nacionais.
\item[Serviços]  Melhor infraestrutura hospitalar e de serviços do departamento (Hospital Regional). Centro administrativo, bancário e educacional (UNVES).
\item[Custo de Vida]  Baixo; economia baseada na pecuária, serviços e agricultura familiar.
\item[Preço da Terra]  US\$ 4.000-7.000/\allowbreak ha (áreas rurais valorizadas); lotes urbanos com preços estáveis de capital departamental.
\item[Hidrologia]  Abundância de mananciais e arroios. Famosa pela fonte Ykua Bolaños. Baixo risco de inundações sistêmicas catastróficas. Drenagem natural favorável.
\item[Clima]  Subtropical Úmido. Verões quentes; invernos secos com geadas ocasionais.
\item[Inclinação recomendada para placas solares]  26° voltado para o Norte (anual); 36° no inverno (Jun-Ago); 16° no verão (Nov-Jan). Base: latitude local -26.18°S.
\item[Energia]  Rede nacional (Itaipu); nodo central de distribuição departamental com alta prioridade de manutenção.
\item[Água]  Abundante via poços artesianos e nascentes históricas (Ykua Bolaños); lençol freático produtivo.
\item[Qualidade do Solo]  Latossolo Vermelho e solos podzólicos; muito férteis e aptos para pecuária de elite, grãos e citros.
\item[Recursos Locais]  Grande excedente de proteína animal (gado de corte), mandioca e frutas. Elevado potencial para autossuficiência hídrica e alimentar básica.
\item[Segurança]  Uma das capitais mais pacíficas do Paraguai. Taxa de homicídios regional entre as menores do país (\textasciitilde 7,0/\allowbreak 100k). Forte segurança institucional e controle social comunitário. Ambiente urbano ordenado e cultural.
\item[Leis Local: Sede de todos os poderes estatais regionais. Livre de Restrição de Fronteira]  Localizada no interior do país, permitindo plena titularidade para brasileiros.
\end{description}
\textbf{Fonte climática:} NASA POWER Climatology API (período 2001-2020)

\textbf{Fonte luminosa:} estimativa world atlas

\paragraph{Irradiação Solar
(kWh/\allowbreak m²/\allowbreak dia)}\label{irradiauxe7uxe3o-solar-kwhmuxb2dia}


\begin{center}
\resizebox{\linewidth}{!}{%
{\renewcommand{\arraystretch}{1.2}%
\begin{tabular}{@{}c|c|c|c|c|c|c|c|c|c|c|c|c@{}}
\toprule
Jan & Fev & Mar & Abr & Mai & Jun & Jul & Ago & Set & Out & Nov & Dez & Média \\
\midrule
6.71 & 6.20 & 5.39 & 4.34 & 3.26 & 2.77 & 3.10 & 3.79 & 4.42 & 5.34 &
6.39 & 6.76 & \textbf{4.87} \\
\bottomrule
\end{tabular}}%
}
\end{center}


\textbf{Inclinação solar recomendada:} 26° N (anual) · 36° N (inverno
jun-ago) · 16° N (verão nov-jan)

\paragraph{Precipitação (mm/\allowbreak mês)}\label{precipitauxe7uxe3o-mmmuxeas}


\begin{center}
\resizebox{\linewidth}{!}{%
{\renewcommand{\arraystretch}{1.2}%
\begin{tabular}{@{}c|c|c|c|c|c|c|c|c|c|c|c|c@{}}
\toprule
Jan & Fev & Mar & Abr & Mai & Jun & Jul & Ago & Set & Out & Nov & Dez & Total/\allowbreak ano \\
\midrule
135 & 129 & 133 & 159 & 166 & 94 & 78 & 61 & 93 & 189 & 188 & 164 &
\textbf{1590 mm} \\
\bottomrule
\end{tabular}}%
}
\end{center}


\paragraph{Poluição Luminosa}\label{poluiuxe7uxe3o-luminosa}

\begin{longtable}[]{@{}ll@{}}
\toprule\noalign{}
Parâmetro & Valor \\
\midrule\noalign{}
\endhead
\bottomrule\noalign{}
\endlastfoot
Escala Bortle & 4 --- Céu rural-suburbano \\
Radiância artificial & 1.6 nW/\allowbreak cm²/\allowbreak sr \\
\end{longtable}
\paragraph{Solo (SoilGrids 2.0, média ponderada 0--30
cm)}

\begin{longtable}[]{@{}ll@{}}
\toprule\noalign{}
Parâmetro & Valor \\
\midrule\noalign{}
\endhead
\bottomrule\noalign{}
\endlastfoot
pH (H$_2$O) & 5.35 \\
Carbono orgânico (SOC) & 16.22 g/\allowbreak kg \\
Argila & 26.0 \% \\
Areia & 46.33 \% \\
Silte (calc.) & 27.7 \% \\
Densidade aparente & 1.35 g/\allowbreak cm³ \\
\textbf{Aptidão agrícola} & \textbf{Média} \\
\end{longtable}

\subsubsection{Indicadores Sociais}

\textbf{Fontes:} PNUD 2020, INE\footnote{\fineINE} Censo 2022, INE\footnote{\fineINE} EPHC 2023, Ministerio
Público 2024\\
\textbf{Nota:} Valores marcados como \emph{(dept.)} referem-se ao
departamento de Caazapá; valores \emph{(dist.)} são específicos deste
distrito. Dados marcados \emph{(est.)} são estimativas.

\begin{longtable}[]{@{}
  >{\raggedright\arraybackslash}p{(\linewidth - 4\tabcolsep) * \real{0.4231}}
  >{\raggedright\arraybackslash}p{(\linewidth - 4\tabcolsep) * \real{0.2692}}
  >{\raggedright\arraybackslash}p{(\linewidth - 4\tabcolsep) * \real{0.3077}}@{}}
\toprule\noalign{}
\begin{minipage}[b]{\linewidth}\raggedright
Indicador
\end{minipage} & \begin{minipage}[b]{\linewidth}\raggedright
Valor
\end{minipage} & \begin{minipage}[b]{\linewidth}\raggedright
Âmbito
\end{minipage} \\
\midrule\noalign{}
\endhead
\bottomrule\noalign{}
\endlastfoot
IDH (2020) & 0,648 & dept. \\
Ranking IDH nacional & 18/\allowbreak 18 & dept. \\
Esperança de vida & 69,5 anos & dept. \\
Escolaridade média & 9.9 anos & dist. \\
RNB per capita (USD PPA) & 5.200 & dept. \\
Pobreza monetária (\%) & 34,7\% & dept. \\
Pobreza extrema (\%) & 11,3\% & dept. \\
Índice de Gini & 0,472 & dept. \\
Acesso a água potável (\%) & 60,1\% & dept. \\
Acesso a saneamento (\%) & 57,4\% & dept. \\
Taxa de homicídios (est., /\allowbreak 100k hab) & \textasciitilde4--6 (estimativa,
próxima à média nacional) & dept. \\
Índice de segurança & médio & dept. \\
População (Censo 2022) & 23.654 hab. & dist. \\
Idade mediana & 31.0 anos & dist. \\
Taxa de fecundidade & N/\allowbreak D & dist. \\
\end{longtable}
\subsubsection{Combustível}

\textbf{Referência:} PETROPAR /\allowbreak  postos locais (2024) \textbf{Tipo de
localidade:} Capital departamental

\begin{longtable}[]{@{}lll@{}}
\toprule\noalign{}
Combustível & USD/\allowbreak litro & Gs/\allowbreak litro (aprox.) \\
\midrule\noalign{}
\endhead
\bottomrule\noalign{}
\endlastfoot
Gasolina 93 oct & 0.99 & 7,326 \\
Gasolina 97 oct (premium) & 1.09 & 8,066 \\
Diesel & 0.91 & 6,734 \\
\end{longtable}

\begin{quote}
Preços podem variar ±5\% conforme posto e sazonalidade. Chaco e interior
remoto apresentam maior variação.
\end{quote}

\subsubsection{Cobertura Celular}

\textbf{Fonte:} CONATEL PY /\allowbreak  operadoras (2024)

\begin{longtable}[]{@{}ll@{}}
\toprule\noalign{}
Parâmetro & Valor \\
\midrule\noalign{}
\endhead
\bottomrule\noalign{}
\endlastfoot
Cobertura 4G população (dept.) & 78\% \\
Cobertura 4G área rural & 55\% \\
Melhor operadora & Tigo \\
Qualidade rural & moderada \\
\end{longtable}

\begin{quote}
Para áreas rurais fora do núcleo urbano, recomenda-se chip Tigo como
principal e Personal como backup.
\end{quote}

\subsubsection{Internet}

\textbf{Fonte:} CONATEL /\allowbreak  Speedtest Ookla (2024)

\begin{longtable}[]{@{}ll@{}}
\toprule\noalign{}
Parâmetro & Valor \\
\midrule\noalign{}
\endhead
\bottomrule\noalign{}
\endlastfoot
Velocidade média download & 32 Mbps \\
Domicílios com internet (dept.) & 42\% \\
Tecnologia predominante & rádio \\
Opção rural & Starlink disponível (\textasciitilde USD 44/\allowbreak mês) \\
\end{longtable}

\subsubsection{Mercado Imobiliário e Terra
Rural}

\textbf{Fonte:} INDERT /\allowbreak  Clasificados.com.py (2024)

\begin{longtable}[]{@{}ll@{}}
\toprule\noalign{}
Tipo & Referência \\
\midrule\noalign{}
\endhead
\bottomrule\noalign{}
\endlastfoot
Terra agrícola alta prod. (USD/\allowbreak ha) & 5,500 \\
Imóvel urbano (USD/\allowbreak m²) & 632 \\
Aluguel 2 quartos (USD/\allowbreak mês) & 264 \\
\end{longtable}

\begin{quote}
Valores de referência departamental. Localidades menores podem ter
preços 20--40\% abaixo da capital departamental. \#\#\# Saúde
\end{quote}

\textbf{Fonte:} MSPBS /\allowbreak  IPS Paraguay (2024-2026), consolidação
departamental e proxy local

\begin{longtable}[]{@{}ll@{}}
\toprule\noalign{}
Serviço & Disponibilidade \\
\midrule\noalign{}
\endhead
\bottomrule\noalign{}
\endlastfoot
USF /\allowbreak  Posto de Saúde & sim \\
Hospital Regional & sim \\
IPS (seguro social) & sim \\
Farmácia & sim \\
Distância ao hospital de referência & local \\
\end{longtable}

\textbf{Principais estabelecimentos:} Hospital distrital /\allowbreak  regional de
Caazapa; rede de USF e farmacias locais

\textbf{Observação para imigrantes:} Centro de referencia do
departamento. Acesso a atendimento primario e, em geral, a maior oferta
publica e privada da area. Cobertura privada continua recomendada para
especialidades e urgências de maior complexidade.
\newpage
\secaoDiagnostico{Doctor Moises Bertoni}{\mapaDistrito{0610}}{Doctor Moisés Bertoni é uma localidade que personifica o isolamento
demográfico em uma das regiões mais pacíficas do Paraguai. Sua maior
vantagem é a invisibilidade estratégica, sendo um distrito pequeno, fora
dos grandes eixos rodoviários, e com uma baixíssima carga populacional.
Oferece um santuário de paz social e recursos naturais básicos
abundantes, exigindo do ocupante autonomia logística e de saúde devido à
infraestrutura local limitada.

\subsubsection{Recomendação
Técnica}

Altamente indicado para estabelecimentos de autossuficiência de
pequeno/\allowbreak médio porte ou residências rurais que funcionem como refúgios de
baixa visibilidade. Requer autonomia em termos de logística básica e
saúde. O GSS de 6.7 reflete a segurança do isolamento em contraste com a
limitação institucional local. Referência atual de score: GSS 6.7.

}
\begin{table}[ht!]
\small \begin{tabular}{p{3.0cm}p{0.8cm}p{6.0cm}} \toprule
\textbf{Critério} & \textbf{Nota} & \textbf{Justificativa} \\ \midrule
A - Ameaças Estratégicas & 6.5 & Localização discreta e protegida de rotas principais; excelente invisibilidade tática \\
B - Risco Social & 7.5 & Baixíssima densidade demográfica; um dos distritos mais vazios da região \\
C - Riscos Naturais & 6.5 & Geologia muito estável; penalizado pelo risco moderado de isolamento pluvial sazonal \\
D - Autossuficiência & 7.0 & Excelentes recursos naturais - gado e água - mas limitado pela infraestrutura logística de insumos \\
E - Institucional & 5.5 & Ambiente institucional seguro e pacífico, mas com infraestrutura pública local limitada \\
\midrule
 \textbf{GSS FINAL} & \textbf{6.7} & \textbf{Moderadamente Seguro (Ideal para quem busca o máximo de isolamento demográfico em uma região segura e produtiva).} \\ \bottomrule
 \end{tabular}
\end{table}
\subsubsection{Vulnerabilidades e
Mitigação}\label{vuln-Doctor_Moises_Bertoni-vulnerabilidades-e-mitigauxe7uxe3o}

\begin{itemize}
\tightlist
\item
  \textbf{Isolamento de Saúde:} Distância significativa de hospitais
  cirúrgicos (Mitigação: manutenção de estoque médico básico e
  treinamento de primeiros socorros avançado).
\item
  \textbf{Dependência Logística:} Necessidade de deslocamento para
  Caazapá para insumos técnicos e comerciais (Mitigação: estoque de
  suprimentos para 60 dias).
\item
  \textbf{Instabilidade Elétrica Rural:} Quedas ocasionais em temporais
  (Mitigação: instalação de sistemas solares fotovoltaicos redundantes).
\end{itemize}

\subsubsection{Dossiê de Campo}\label{dossie-Doctor_Moises_Bertoni-dossiuxea-de-campo}
\begin{description}[style=nextline,leftmargin=0.5cm,font=\bfseries]
\item[Coordenadas]  26°28'S 56°20'W.
\item[Topografia]  Relevo predominantemente plano a suavemente ondulado, localizado nas proximidades da bacia do Rio Tebicuary. Altitude \textasciitilde 140m. Área de \textasciitilde 637 km². Geologicamente muito estável, risco sísmico desprezível.
\item[Alvos Estratégicos]  Áreas de produção pecuária extensiva; Pequenas colônias de agricultura tradicional. Ausência absoluta de alvos militares ou industriais primários, oferecendo altíssima invisibilidade tática e discrição estratégica.
\item[Fallout]  Ventos predominantes do Norte (Viento Norte, quente/\allowbreak seco no verão) e Leste. Baixo risco de fallout direto.
\item[População]  5.960 habitantes (Censo 2022 Final). População em declínio, com perfil essencialmente rural (82\%) e laborioso.
\item[Densidade]  \textasciitilde 9,4 hab/\allowbreak km² (Densidade baixíssima, oferecendo isolamento social superior e privacidade estratégica absoluta fora dos núcleos povoados).
\item[Vias de Acesso]  Acesso via ramal asfaltado de 13 km conectando à Ruta PY08 (Blas Garay). Conectividade interna dependente de estradas vicinais de terra que podem sofrer isolamento parcial em períodos de chuva severa (média pluvial 1.600 mm/\allowbreak ano).
\item[Serviços]  Unidade de Saúde da Família (USF) local para atendimentos básicos. Dependência crítica da capital Caazapá (25 km) para serviços médicos de alta complexidade e logística comercial diversificada.
\item[Custo de Vida]  Muito baixo; economia baseada na pecuária de subsistência e pequena agricultura.
\item[Preço da Terra]  US\$ 2.500-4.000/\allowbreak ha (áreas de pecuária extensiva); US\$ 5.000/\allowbreak ha (áreas mecanizáveis ou próximas ao asfalto).
\item[Hidrologia]  Baixo risco de inundações sistêmicas catastróficas. Presença de arroios locais tributários do Rio Tebicuary que favorecem a pecuária. Risco de erosão em caminhos vicinais em eventos climáticos extremos.
\item[Clima]  Subtropical Úmido. Verões quentes; invernos secos com entradas bruscas de ar frio do Sul.
\item[Energia]  Rede nacional (Itaipu); instável em áreas rurais devido ao isolamento da rede.
\item[Água]  Abundante via poços artesianos e lençol freático superficial; excelente qualidade hídrica subterrânea.
\item[Qualidade do Solo]  Solo franco-arenoso a argiloso (Podzólicos); fértil e profundo, ideal para pastagens, mandioca e milho.
\item[Recursos Locais]  Produção de gado de corte, leite e alimentos básicos de subsistência. Elevadíssimo potencial para autossuficiência alimentar básica per capita devido à baixíssima carga demográfica.
\item[Segurança]  Zona rural extremamente pacífica e segura. Taxa de homicídios no departamento de Caazapá entre as menores do Paraguai (\textasciitilde 7,0/\allowbreak 100k). Coesão social forte baseada em comunidades tradicionais. Baixíssima visibilidade para conflitos externos.
\item[Leis Local: Município histórico e consolidado. Livre de Restrição de Fronteira]  Localizado no interior do país, permitindo plena titularidade para investidores brasileiros.
\end{description}
\textbf{Fonte climática:} NASA POWER Climatology API (período 2001-2020)

\textbf{Fonte luminosa:} estimativa world atlas

\paragraph{Irradiação Solar
(kWh/\allowbreak m²/\allowbreak dia)}\label{irradiauxe7uxe3o-solar-kwhmuxb2dia}


\begin{center}
\resizebox{\linewidth}{!}{%
{\renewcommand{\arraystretch}{1.2}%
\begin{tabular}{@{}c|c|c|c|c|c|c|c|c|c|c|c|c@{}}
\toprule
Jan & Fev & Mar & Abr & Mai & Jun & Jul & Ago & Set & Out & Nov & Dez & Média \\
\midrule
6.71 & 6.20 & 5.39 & 4.34 & 3.26 & 2.77 & 3.10 & 3.79 & 4.42 & 5.34 &
6.39 & 6.76 & \textbf{4.87} \\
\bottomrule
\end{tabular}}%
}
\end{center}


\textbf{Inclinação solar recomendada:} 26° N (anual) · 36° N (inverno
jun-ago) · 16° N (verão nov-jan)

\paragraph{Precipitação (mm/\allowbreak mês)}\label{precipitauxe7uxe3o-mmmuxeas}


\begin{center}
\resizebox{\linewidth}{!}{%
{\renewcommand{\arraystretch}{1.2}%
\begin{tabular}{@{}c|c|c|c|c|c|c|c|c|c|c|c|c@{}}
\toprule
Jan & Fev & Mar & Abr & Mai & Jun & Jul & Ago & Set & Out & Nov & Dez & Total/\allowbreak ano \\
\midrule
141 & 123 & 136 & 171 & 156 & 100 & 82 & 71 & 100 & 194 & 188 & 167 &
\textbf{1629 mm} \\
\bottomrule
\end{tabular}}%
}
\end{center}


\paragraph{Poluição Luminosa}\label{poluiuxe7uxe3o-luminosa}

\begin{longtable}[]{@{}ll@{}}
\toprule\noalign{}
Parâmetro & Valor \\
\midrule\noalign{}
\endhead
\bottomrule\noalign{}
\endlastfoot
Escala Bortle & 3 --- Céu rural \\
Radiância artificial & 0.8 nW/\allowbreak cm²/\allowbreak sr \\
\end{longtable}
\paragraph{Solo (SoilGrids 2.0, média ponderada 0--30
cm)}

\begin{longtable}[]{@{}ll@{}}
\toprule\noalign{}
Parâmetro & Valor \\
\midrule\noalign{}
\endhead
\bottomrule\noalign{}
\endlastfoot
pH (H$_2$O) & 5.2 \\
Carbono orgânico (SOC) & 16.55 g/\allowbreak kg \\
Argila & 21.52 \% \\
Areia & 51.66 \% \\
Silte (calc.) & 26.8 \% \\
Densidade aparente & 1.37 g/\allowbreak cm³ \\
\textbf{Aptidão agrícola} & \textbf{Média} \\
\end{longtable}

\subsubsection{Indicadores Sociais}

\textbf{Fontes:} PNUD 2020, INE\footnote{\fineINE} Censo 2022, INE\footnote{\fineINE} EPHC 2023, Ministerio
Público 2024\\
\textbf{Nota:} Valores marcados como \emph{(dept.)} referem-se ao
departamento de Caazapá; valores \emph{(dist.)} são específicos deste
distrito. Dados marcados \emph{(est.)} são estimativas.

\begin{longtable}[]{@{}
  >{\raggedright\arraybackslash}p{(\linewidth - 4\tabcolsep) * \real{0.4231}}
  >{\raggedright\arraybackslash}p{(\linewidth - 4\tabcolsep) * \real{0.2692}}
  >{\raggedright\arraybackslash}p{(\linewidth - 4\tabcolsep) * \real{0.3077}}@{}}
\toprule\noalign{}
\begin{minipage}[b]{\linewidth}\raggedright
Indicador
\end{minipage} & \begin{minipage}[b]{\linewidth}\raggedright
Valor
\end{minipage} & \begin{minipage}[b]{\linewidth}\raggedright
Âmbito
\end{minipage} \\
\midrule\noalign{}
\endhead
\bottomrule\noalign{}
\endlastfoot
IDH (2020) & 0,648 & dept. \\
Ranking IDH nacional & 18/\allowbreak 18 & dept. \\
Esperança de vida & 69,5 anos & dept. \\
Escolaridade média & 6,4 anos & dept. \\
RNB per capita (USD PPA) & 5.200 & dept. \\
Pobreza monetária (\%) & 34,7\% & dept. \\
Pobreza extrema (\%) & 11,3\% & dept. \\
Índice de Gini & 0,472 & dept. \\
Acesso a água potável (\%) & 60,1\% & dept. \\
Acesso a saneamento (\%) & 57,4\% & dept. \\
Taxa de homicídios (est., /\allowbreak 100k hab) & \textasciitilde4--6 (estimativa,
próxima à média nacional) & dept. \\
Índice de segurança & médio & dept. \\
População (Censo 2022) & 3.815 hab. & dist. \\
Idade mediana & 29.0 anos & dist. \\
Taxa de fecundidade & N/\allowbreak D & dist. \\
\end{longtable}
\subsubsection{Combustível}

\textbf{Referência:} PETROPAR /\allowbreak  postos locais (2024) \textbf{Tipo de
localidade:} Interior

\begin{longtable}[]{@{}lll@{}}
\toprule\noalign{}
Combustível & USD/\allowbreak litro & Gs/\allowbreak litro (aprox.) \\
\midrule\noalign{}
\endhead
\bottomrule\noalign{}
\endlastfoot
Gasolina 93 oct & 0.99 & 7,326 \\
Gasolina 97 oct (premium) & 1.09 & 8,066 \\
Diesel & 0.91 & 6,734 \\
\end{longtable}

\begin{quote}
Preços podem variar ±5\% conforme posto e sazonalidade. Chaco e interior
remoto apresentam maior variação.
\end{quote}

\subsubsection{Cobertura Celular}

\textbf{Fonte:} CONATEL PY /\allowbreak  operadoras (2024)

\begin{longtable}[]{@{}ll@{}}
\toprule\noalign{}
Parâmetro & Valor \\
\midrule\noalign{}
\endhead
\bottomrule\noalign{}
\endlastfoot
Cobertura 4G população (dept.) & 78\% \\
Cobertura 4G área rural & 55\% \\
Melhor operadora & Tigo \\
Qualidade rural & moderada \\
\end{longtable}

\begin{quote}
Para áreas rurais fora do núcleo urbano, recomenda-se chip Tigo como
principal e Personal como backup.
\end{quote}

\subsubsection{Internet}

\textbf{Fonte:} CONATEL /\allowbreak  Speedtest Ookla (2024)

\begin{longtable}[]{@{}ll@{}}
\toprule\noalign{}
Parâmetro & Valor \\
\midrule\noalign{}
\endhead
\bottomrule\noalign{}
\endlastfoot
Velocidade média download & 32 Mbps \\
Domicílios com internet (dept.) & 42\% \\
Tecnologia predominante & rádio \\
Opção rural & Starlink disponível (\textasciitilde USD 44/\allowbreak mês) \\
\end{longtable}

\subsubsection{Mercado Imobiliário e Terra
Rural}

\textbf{Fonte:} INDERT /\allowbreak  Clasificados.com.py (2024)

\begin{longtable}[]{@{}ll@{}}
\toprule\noalign{}
Tipo & Referência \\
\midrule\noalign{}
\endhead
\bottomrule\noalign{}
\endlastfoot
Terra agrícola alta prod. (USD/\allowbreak ha) & 5,500 \\
Imóvel urbano (USD/\allowbreak m²) & 550 \\
Aluguel 2 quartos (USD/\allowbreak mês) & 220 \\
\end{longtable}

\begin{quote}
Valores de referência departamental. Localidades menores podem ter
preços 20--40\% abaixo da capital departamental. \#\#\# Saúde
\end{quote}

\textbf{Fonte:} MSPBS /\allowbreak  IPS Paraguay (2024-2026), consolidação
departamental e proxy local

\begin{longtable}[]{@{}ll@{}}
\toprule\noalign{}
Serviço & Disponibilidade \\
\midrule\noalign{}
\endhead
\bottomrule\noalign{}
\endlastfoot
USF /\allowbreak  Posto de Saúde & sim \\
Hospital Regional & não \\
IPS (seguro social) & não \\
Farmácia & sim \\
Distância ao hospital de referência & \textasciitilde27 km (Caazapa) \\
\end{longtable}

\textbf{Principais estabelecimentos:} USF local; referencia hospitalar
em Caazapa

\textbf{Observação para imigrantes:} Atendimento primario local ou em
raio curto; casos de maior complexidade seguem para Caazapa. Cobertura
privada continua recomendada para especialidades e urgências de maior
complexidade.
\newpage
\secaoDiagnostico{Fulgencio Yegros}{\mapaDistrito{0609}}{Fulgencio Yegros (Caazapa) apresenta perfil operacional consolidado para
comparacao territorial, com leitura conjunta de infraestrutura, risco
hidroclimatico e capacidade de continuidade de servicos essenciais.

\subsubsection{Pre-condicoes Minimas de
Decisao}

\begin{itemize}
\tightlist
\item
  Atualizar indicadores criticos em ciclo trimestral.
\item
  Confirmar trafegabilidade e contingencia hidrica/\allowbreak energetica local.
\item
  Validar checklist de resiliencia domiciliar/\allowbreak logistica antes de decisao
  final.
\end{itemize}

\subsubsection{Recomendação
Técnica}

Uso recomendado como candidato em analise multicriterio, condicionado ao
cumprimento das pre-condicoes acima. Referência atual de score: GSS 7.7;
risco predominante: moderado.

}
\begin{table}[ht!]
\small \begin{tabular}{p{3.0cm}p{0.8cm}p{6.0cm}} \toprule
\textbf{Critério} & \textbf{Nota} & \textbf{Justificativa} \\ \midrule
A - Ameaças Estratégicas & 7.8 & - \\
B - Risco Social & 8.7 & - \\
C - Riscos Naturais & 6.2 & - \\
D - Autossuficiência & 7.2 & - \\
E - Institucional & 8.4 & - \\
\midrule
 \textbf{GSS FINAL} & \textbf{7.7} & \textbf{Seguro (bom potencial com preparacao).} \\ \bottomrule
 \end{tabular}
\end{table}
\subsubsection{Vulnerabilidades e
Mitigação}\label{vuln-Fulgencio_Yegros-vulnerabilidades-e-mitigauxe7uxe3o}

\begin{itemize}
\tightlist
\item
  Exposicao hidrometeorologica moderada (45-64/\allowbreak 100).
\item
  Conectividade viaria favoravel (\textgreater=70/\allowbreak 100).
\item
  Estoque de contingencia para 14 dias (agua/\allowbreak alimentos/\allowbreak combustivel),
  priorizado quando risco hidro=47/\allowbreak 100.
\item
  Duplicar rotas/\allowbreak logistica quando conectividade viaria=72/\allowbreak 100 for
  \textless70; manter rota secundaria mapeada e testada trimestralmente.
\item
  Plano de energia resiliente com autonomia minima de 72h
  (geracao/\allowbreak backup), revisao semestral.
\end{itemize}

\subsubsection{Dossiê de Campo}\label{dossie-Fulgencio_Yegros-dossiuxea-de-campo}
\begin{description}[style=nextline,leftmargin=0.5cm,font=\bfseries]
\item[Coordenadas]  26°27'S 56°24'W (geocodificado via Nominatim).
\end{description}
\textbf{Fonte climática:} NASA POWER Climatology API (período 2001-2020)

\textbf{Fonte luminosa:} estimativa world atlas

\paragraph{Irradiação Solar
(kWh/\allowbreak m²/\allowbreak dia)}\label{irradiauxe7uxe3o-solar-kwhmuxb2dia}


\begin{center}
\resizebox{\linewidth}{!}{%
{\renewcommand{\arraystretch}{1.2}%
\begin{tabular}{@{}c|c|c|c|c|c|c|c|c|c|c|c|c@{}}
\toprule
Jan & Fev & Mar & Abr & Mai & Jun & Jul & Ago & Set & Out & Nov & Dez & Média \\
\midrule
6.71 & 6.20 & 5.39 & 4.34 & 3.26 & 2.77 & 3.10 & 3.79 & 4.42 & 5.34 &
6.39 & 6.76 & \textbf{4.87} \\
\bottomrule
\end{tabular}}%
}
\end{center}


\textbf{Inclinação solar recomendada:} 26° N (anual) · 36° N (inverno
jun-ago) · 16° N (verão nov-jan)

\paragraph{Precipitação (mm/\allowbreak mês)}\label{precipitauxe7uxe3o-mmmuxeas}


\begin{center}
\resizebox{\linewidth}{!}{%
{\renewcommand{\arraystretch}{1.2}%
\begin{tabular}{@{}c|c|c|c|c|c|c|c|c|c|c|c|c@{}}
\toprule
Jan & Fev & Mar & Abr & Mai & Jun & Jul & Ago & Set & Out & Nov & Dez & Total/\allowbreak ano \\
\midrule
141 & 123 & 136 & 171 & 156 & 100 & 82 & 71 & 100 & 194 & 188 & 167 &
\textbf{1629 mm} \\
\bottomrule
\end{tabular}}%
}
\end{center}


\paragraph{Poluição Luminosa}\label{poluiuxe7uxe3o-luminosa}

\begin{longtable}[]{@{}ll@{}}
\toprule\noalign{}
Parâmetro & Valor \\
\midrule\noalign{}
\endhead
\bottomrule\noalign{}
\endlastfoot
Escala Bortle & 3 --- Céu rural \\
Radiância artificial & 0.8 nW/\allowbreak cm²/\allowbreak sr \\
\end{longtable}
\paragraph{Solo (SoilGrids 2.0, média ponderada 0--30
cm)}

\begin{longtable}[]{@{}ll@{}}
\toprule\noalign{}
Parâmetro & Valor \\
\midrule\noalign{}
\endhead
\bottomrule\noalign{}
\endlastfoot
pH (H$_2$O) & 5.3 \\
Carbono orgânico (SOC) & 15.88 g/\allowbreak kg \\
Argila & 23.73 \% \\
Areia & 48.22 \% \\
Silte (calc.) & 28.0 \% \\
Densidade aparente & 1.39 g/\allowbreak cm³ \\
\textbf{Aptidão agrícola} & \textbf{Média} \\
\end{longtable}

\subsubsection{Indicadores Sociais}

\textbf{Fontes:} PNUD 2020, INE\footnote{\fineINE} Censo 2022, INE\footnote{\fineINE} EPHC 2023, Ministerio
Público 2024\\
\textbf{Nota:} Valores marcados como \emph{(dept.)} referem-se ao
departamento de Caazapá; valores \emph{(dist.)} são específicos deste
distrito. Dados marcados \emph{(est.)} são estimativas.

\begin{longtable}[]{@{}
  >{\raggedright\arraybackslash}p{(\linewidth - 4\tabcolsep) * \real{0.4231}}
  >{\raggedright\arraybackslash}p{(\linewidth - 4\tabcolsep) * \real{0.2692}}
  >{\raggedright\arraybackslash}p{(\linewidth - 4\tabcolsep) * \real{0.3077}}@{}}
\toprule\noalign{}
\begin{minipage}[b]{\linewidth}\raggedright
Indicador
\end{minipage} & \begin{minipage}[b]{\linewidth}\raggedright
Valor
\end{minipage} & \begin{minipage}[b]{\linewidth}\raggedright
Âmbito
\end{minipage} \\
\midrule\noalign{}
\endhead
\bottomrule\noalign{}
\endlastfoot
IDH (2020) & 0,648 & dept. \\
Ranking IDH nacional & 18/\allowbreak 18 & dept. \\
Esperança de vida & 69,5 anos & dept. \\
Escolaridade média & 6,4 anos & dept. \\
RNB per capita (USD PPA) & 5.200 & dept. \\
Pobreza monetária (\%) & 34,7\% & dept. \\
Pobreza extrema (\%) & 11,3\% & dept. \\
Índice de Gini & 0,472 & dept. \\
Acesso a água potável (\%) & 60,1\% & dept. \\
Acesso a saneamento (\%) & 57,4\% & dept. \\
Taxa de homicídios (est., /\allowbreak 100k hab) & \textasciitilde4--6 (estimativa,
próxima à média nacional) & dept. \\
Índice de segurança & médio & dept. \\
População (Censo 2022) & N/\allowbreak D & dist. \\
Idade mediana & N/\allowbreak D & dist. \\
Taxa de fecundidade & N/\allowbreak D & dist. \\
\end{longtable}
\subsubsection{Combustível}

\textbf{Referência:} PETROPAR /\allowbreak  postos locais (2024) \textbf{Tipo de
localidade:} Interior

\begin{longtable}[]{@{}lll@{}}
\toprule\noalign{}
Combustível & USD/\allowbreak litro & Gs/\allowbreak litro (aprox.) \\
\midrule\noalign{}
\endhead
\bottomrule\noalign{}
\endlastfoot
Gasolina 93 oct & 0.99 & 7,326 \\
Gasolina 97 oct (premium) & 1.09 & 8,066 \\
Diesel & 0.91 & 6,734 \\
\end{longtable}

\begin{quote}
Preços podem variar ±5\% conforme posto e sazonalidade. Chaco e interior
remoto apresentam maior variação.
\end{quote}

\subsubsection{Cobertura Celular}

\textbf{Fonte:} CONATEL PY /\allowbreak  operadoras (2024)

\begin{longtable}[]{@{}ll@{}}
\toprule\noalign{}
Parâmetro & Valor \\
\midrule\noalign{}
\endhead
\bottomrule\noalign{}
\endlastfoot
Cobertura 4G população (dept.) & 78\% \\
Cobertura 4G área rural & 55\% \\
Melhor operadora & Tigo \\
Qualidade rural & moderada \\
\end{longtable}

\begin{quote}
Para áreas rurais fora do núcleo urbano, recomenda-se chip Tigo como
principal e Personal como backup.
\end{quote}

\subsubsection{Internet}

\textbf{Fonte:} CONATEL /\allowbreak  Speedtest Ookla (2024)

\begin{longtable}[]{@{}ll@{}}
\toprule\noalign{}
Parâmetro & Valor \\
\midrule\noalign{}
\endhead
\bottomrule\noalign{}
\endlastfoot
Velocidade média download & 32 Mbps \\
Domicílios com internet (dept.) & 42\% \\
Tecnologia predominante & rádio \\
Opção rural & Starlink disponível (\textasciitilde USD 44/\allowbreak mês) \\
\end{longtable}

\subsubsection{Mercado Imobiliário e Terra
Rural}

\textbf{Fonte:} INDERT /\allowbreak  Clasificados.com.py (2024)

\begin{longtable}[]{@{}ll@{}}
\toprule\noalign{}
Tipo & Referência \\
\midrule\noalign{}
\endhead
\bottomrule\noalign{}
\endlastfoot
Terra agrícola alta prod. (USD/\allowbreak ha) & 5,500 \\
Imóvel urbano (USD/\allowbreak m²) & 550 \\
Aluguel 2 quartos (USD/\allowbreak mês) & 220 \\
\end{longtable}

\begin{quote}
Valores de referência departamental. Localidades menores podem ter
preços 20--40\% abaixo da capital departamental. \#\#\# Saúde
\end{quote}

\textbf{Fonte:} MSPBS /\allowbreak  IPS Paraguay (2024-2026), consolidação
departamental e proxy local

\begin{longtable}[]{@{}ll@{}}
\toprule\noalign{}
Serviço & Disponibilidade \\
\midrule\noalign{}
\endhead
\bottomrule\noalign{}
\endlastfoot
USF /\allowbreak  Posto de Saúde & sim \\
Hospital Regional & não \\
IPS (seguro social) & não \\
Farmácia & sim \\
Distância ao hospital de referência & \textasciitilde36 km (Caazapa) \\
\end{longtable}

\textbf{Principais estabelecimentos:} USF local; referencia hospitalar
em Caazapa

\textbf{Observação para imigrantes:} Atendimento primario local ou em
raio curto; casos de maior complexidade seguem para Caazapa. Cobertura
privada continua recomendada para especialidades e urgências de maior
complexidade.
\newpage
\secaoDiagnostico{General Higinio Morinigo}{\mapaDistrito{0605}}{General Higinio Morinigo (Caazapa) apresenta perfil operacional
consolidado para comparacao territorial, com leitura conjunta de
infraestrutura, risco hidroclimatico e capacidade de continuidade de
servicos essenciais.

\subsubsection{Pre-condicoes Minimas de
Decisao}

\begin{itemize}
\tightlist
\item
  Atualizar indicadores criticos em ciclo trimestral.
\item
  Confirmar trafegabilidade e contingencia hidrica/\allowbreak energetica local.
\item
  Validar checklist de resiliencia domiciliar/\allowbreak logistica antes de decisao
  final.
\end{itemize}

\subsubsection{Recomendação
Técnica}

Uso recomendado como candidato em analise multicriterio, condicionado ao
cumprimento das pre-condicoes acima. Referência atual de score: GSS 7.1;
risco predominante: baixo.

}
\begin{table}[ht!]
\small \begin{tabular}{p{3.0cm}p{0.8cm}p{6.0cm}} \toprule
\textbf{Critério} & \textbf{Nota} & \textbf{Justificativa} \\ \midrule
A - Ameaças Estratégicas & 7.4 & - \\
B - Risco Social & 6.7 & - \\
C - Riscos Naturais & 7.1 & - \\
D - Autossuficiência & 6.7 & - \\
E - Institucional & 7.4 & - \\
\midrule
 \textbf{GSS FINAL} & \textbf{7.1} & \textbf{Seguro (bom potencial com preparacao).} \\ \bottomrule
 \end{tabular}
\end{table}
\subsubsection{Vulnerabilidades e
Mitigação}\label{vuln-General_Higinio_Morinigo-vulnerabilidades-e-mitigauxe7uxe3o}

\begin{itemize}
\tightlist
\item
  Exposicao hidrometeorologica controlada (\textless45/\allowbreak 100).
\item
  Conectividade viaria intermediaria (50-69/\allowbreak 100).
\item
  Estoque de contingencia para 14 dias (agua/\allowbreak alimentos/\allowbreak combustivel),
  priorizado quando risco hidro=36/\allowbreak 100.
\item
  Duplicar rotas/\allowbreak logistica quando conectividade viaria=50/\allowbreak 100 for
  \textless70; manter rota secundaria mapeada e testada trimestralmente.
\item
  Plano de energia resiliente com autonomia minima de 72h
  (geracao/\allowbreak backup), revisao semestral.
\end{itemize}

\subsubsection{Dossiê de Campo}\label{dossie-General_Higinio_Morinigo-dossiuxea-de-campo}
\begin{description}[style=nextline,leftmargin=0.5cm,font=\bfseries]
\item[Coordenadas]  26°02'S 56°05'W (geocodificado via Nominatim).
\end{description}
\textbf{Fonte climática:} NASA POWER Climatology API (período 2001-2020)

\textbf{Fonte luminosa:} estimativa world atlas

\paragraph{Irradiação Solar
(kWh/\allowbreak m²/\allowbreak dia)}\label{irradiauxe7uxe3o-solar-kwhmuxb2dia}


\begin{center}
\resizebox{\linewidth}{!}{%
{\renewcommand{\arraystretch}{1.2}%
\begin{tabular}{@{}c|c|c|c|c|c|c|c|c|c|c|c|c@{}}
\toprule
Jan & Fev & Mar & Abr & Mai & Jun & Jul & Ago & Set & Out & Nov & Dez & Média \\
\midrule
6.71 & 6.20 & 5.39 & 4.34 & 3.26 & 2.77 & 3.10 & 3.79 & 4.42 & 5.34 &
6.39 & 6.76 & \textbf{4.87} \\
\bottomrule
\end{tabular}}%
}
\end{center}


\textbf{Inclinação solar recomendada:} 26° N (anual) · 36° N (inverno
jun-ago) · 16° N (verão nov-jan)

\paragraph{Precipitação (mm/\allowbreak mês)}\label{precipitauxe7uxe3o-mmmuxeas}


\begin{center}
\resizebox{\linewidth}{!}{%
{\renewcommand{\arraystretch}{1.2}%
\begin{tabular}{@{}c|c|c|c|c|c|c|c|c|c|c|c|c@{}}
\toprule
Jan & Fev & Mar & Abr & Mai & Jun & Jul & Ago & Set & Out & Nov & Dez & Total/\allowbreak ano \\
\midrule
135 & 129 & 133 & 159 & 166 & 94 & 78 & 61 & 93 & 189 & 188 & 164 &
\textbf{1590 mm} \\
\bottomrule
\end{tabular}}%
}
\end{center}


\paragraph{Poluição Luminosa}\label{poluiuxe7uxe3o-luminosa}

\begin{longtable}[]{@{}ll@{}}
\toprule\noalign{}
Parâmetro & Valor \\
\midrule\noalign{}
\endhead
\bottomrule\noalign{}
\endlastfoot
Escala Bortle & 3 --- Céu rural \\
Radiância artificial & 0.8 nW/\allowbreak cm²/\allowbreak sr \\
\end{longtable}
\paragraph{Solo (SoilGrids 2.0, média ponderada 0--30
cm)}

\begin{longtable}[]{@{}ll@{}}
\toprule\noalign{}
Parâmetro & Valor \\
\midrule\noalign{}
\endhead
\bottomrule\noalign{}
\endlastfoot
pH (H$_2$O) & 5.3 \\
Carbono orgânico (SOC) & 17.99 g/\allowbreak kg \\
Argila & 29.18 \% \\
Areia & 39.96 \% \\
Silte (calc.) & 30.9 \% \\
Densidade aparente & 1.32 g/\allowbreak cm³ \\
\textbf{Aptidão agrícola} & \textbf{Média} \\
\end{longtable}

\subsubsection{Indicadores Sociais}

\textbf{Fontes:} PNUD 2020, INE\footnote{\fineINE} Censo 2022, INE\footnote{\fineINE} EPHC 2023, Ministerio
Público 2024\\
\textbf{Nota:} Valores marcados como \emph{(dept.)} referem-se ao
departamento de Caazapá; valores \emph{(dist.)} são específicos deste
distrito. Dados marcados \emph{(est.)} são estimativas.

\begin{longtable}[]{@{}
  >{\raggedright\arraybackslash}p{(\linewidth - 4\tabcolsep) * \real{0.4231}}
  >{\raggedright\arraybackslash}p{(\linewidth - 4\tabcolsep) * \real{0.2692}}
  >{\raggedright\arraybackslash}p{(\linewidth - 4\tabcolsep) * \real{0.3077}}@{}}
\toprule\noalign{}
\begin{minipage}[b]{\linewidth}\raggedright
Indicador
\end{minipage} & \begin{minipage}[b]{\linewidth}\raggedright
Valor
\end{minipage} & \begin{minipage}[b]{\linewidth}\raggedright
Âmbito
\end{minipage} \\
\midrule\noalign{}
\endhead
\bottomrule\noalign{}
\endlastfoot
IDH (2020) & 0,648 & dept. \\
Ranking IDH nacional & 18/\allowbreak 18 & dept. \\
Esperança de vida & 69,5 anos & dept. \\
Escolaridade média & 6,4 anos & dept. \\
RNB per capita (USD PPA) & 5.200 & dept. \\
Pobreza monetária (\%) & 34,7\% & dept. \\
Pobreza extrema (\%) & 11,3\% & dept. \\
Índice de Gini & 0,472 & dept. \\
Acesso a água potável (\%) & 60,1\% & dept. \\
Acesso a saneamento (\%) & 57,4\% & dept. \\
Taxa de homicídios (est., /\allowbreak 100k hab) & \textasciitilde4--6 (estimativa,
próxima à média nacional) & dept. \\
Índice de segurança & médio & dept. \\
População (Censo 2022) & 4.728 hab. & dist. \\
Idade mediana & 32.0 anos & dist. \\
Taxa de fecundidade & N/\allowbreak D & dist. \\
\end{longtable}
\subsubsection{Combustível}

\textbf{Referência:} PETROPAR /\allowbreak  postos locais (2024) \textbf{Tipo de
localidade:} Interior

\begin{longtable}[]{@{}lll@{}}
\toprule\noalign{}
Combustível & USD/\allowbreak litro & Gs/\allowbreak litro (aprox.) \\
\midrule\noalign{}
\endhead
\bottomrule\noalign{}
\endlastfoot
Gasolina 93 oct & 0.99 & 7,326 \\
Gasolina 97 oct (premium) & 1.09 & 8,066 \\
Diesel & 0.91 & 6,734 \\
\end{longtable}

\begin{quote}
Preços podem variar ±5\% conforme posto e sazonalidade. Chaco e interior
remoto apresentam maior variação.
\end{quote}

\subsubsection{Cobertura Celular}

\textbf{Fonte:} CONATEL PY /\allowbreak  operadoras (2024)

\begin{longtable}[]{@{}ll@{}}
\toprule\noalign{}
Parâmetro & Valor \\
\midrule\noalign{}
\endhead
\bottomrule\noalign{}
\endlastfoot
Cobertura 4G população (dept.) & 78\% \\
Cobertura 4G área rural & 55\% \\
Melhor operadora & Tigo \\
Qualidade rural & moderada \\
\end{longtable}

\begin{quote}
Para áreas rurais fora do núcleo urbano, recomenda-se chip Tigo como
principal e Personal como backup.
\end{quote}

\subsubsection{Internet}

\textbf{Fonte:} CONATEL /\allowbreak  Speedtest Ookla (2024)

\begin{longtable}[]{@{}ll@{}}
\toprule\noalign{}
Parâmetro & Valor \\
\midrule\noalign{}
\endhead
\bottomrule\noalign{}
\endlastfoot
Velocidade média download & 32 Mbps \\
Domicílios com internet (dept.) & 42\% \\
Tecnologia predominante & rádio \\
Opção rural & Starlink disponível (\textasciitilde USD 44/\allowbreak mês) \\
\end{longtable}

\subsubsection{Mercado Imobiliário e Terra
Rural}

\textbf{Fonte:} INDERT /\allowbreak  Clasificados.com.py (2024)

\begin{longtable}[]{@{}ll@{}}
\toprule\noalign{}
Tipo & Referência \\
\midrule\noalign{}
\endhead
\bottomrule\noalign{}
\endlastfoot
Terra agrícola alta prod. (USD/\allowbreak ha) & 5,500 \\
Imóvel urbano (USD/\allowbreak m²) & 550 \\
Aluguel 2 quartos (USD/\allowbreak mês) & 220 \\
\end{longtable}

\begin{quote}
Valores de referência departamental. Localidades menores podem ter
preços 20--40\% abaixo da capital departamental. \#\#\# Saúde
\end{quote}

\textbf{Fonte:} MSPBS /\allowbreak  IPS Paraguay (2024-2026), consolidação
departamental e proxy local

\begin{longtable}[]{@{}ll@{}}
\toprule\noalign{}
Serviço & Disponibilidade \\
\midrule\noalign{}
\endhead
\bottomrule\noalign{}
\endlastfoot
USF /\allowbreak  Posto de Saúde & sim \\
Hospital Regional & não \\
IPS (seguro social) & não \\
Farmácia & sim \\
Distância ao hospital de referência & \textasciitilde42 km (Caazapa) \\
\end{longtable}

\textbf{Principais estabelecimentos:} USF local; referencia hospitalar
em Caazapa

\textbf{Observação para imigrantes:} Atendimento primario local ou em
raio curto; casos de maior complexidade seguem para Caazapa. Cobertura
privada continua recomendada para especialidades e urgências de maior
complexidade.
\newpage
\secaoDiagnostico{Maciel}{\mapaDistrito{0606}}{Maciel (Caazapa) apresenta perfil operacional consolidado para
comparacao territorial, com leitura conjunta de infraestrutura, risco
hidroclimatico e capacidade de continuidade de servicos essenciais.

\subsubsection{Pre-condicoes Minimas de
Decisao}

\begin{itemize}
\tightlist
\item
  Atualizar indicadores criticos em ciclo trimestral.
\item
  Confirmar trafegabilidade e contingencia hidrica/\allowbreak energetica local.
\item
  Validar checklist de resiliencia domiciliar/\allowbreak logistica antes de decisao
  final.
\end{itemize}

\subsubsection{Recomendação
Técnica}

Uso recomendado como candidato em analise multicriterio, condicionado ao
cumprimento das pre-condicoes acima. Referência atual de score: GSS 7.4;
risco predominante: moderado.

}
\begin{table}[ht!]
\small \begin{tabular}{p{3.0cm}p{0.8cm}p{6.0cm}} \toprule
\textbf{Critério} & \textbf{Nota} & \textbf{Justificativa} \\ \midrule
A - Ameaças Estratégicas & 7.6 & - \\
B - Risco Social & 8.7 & - \\
C - Riscos Naturais & 5.3 & - \\
D - Autossuficiência & 6.5 & - \\
E - Institucional & 8.1 & - \\
\midrule
 \textbf{GSS FINAL} & \textbf{7.4} & \textbf{Seguro (bom potencial com preparacao).} \\ \bottomrule
 \end{tabular}
\end{table}
\subsubsection{Vulnerabilidades e
Mitigação}\label{vuln-Maciel-vulnerabilidades-e-mitigauxe7uxe3o}

\begin{itemize}
\tightlist
\item
  Exposicao hidrometeorologica moderada (45-64/\allowbreak 100).
\item
  Conectividade viaria favoravel (\textgreater=70/\allowbreak 100).
\item
  Estoque de contingencia para 14 dias (agua/\allowbreak alimentos/\allowbreak combustivel),
  priorizado quando risco hidro=59/\allowbreak 100.
\item
  Duplicar rotas/\allowbreak logistica quando conectividade viaria=80/\allowbreak 100 for
  \textless70; manter rota secundaria mapeada e testada trimestralmente.
\item
  Plano de energia resiliente com autonomia minima de 72h
  (geracao/\allowbreak backup), revisao semestral.
\end{itemize}

\subsubsection{Dossiê de Campo}\label{dossie-Maciel-dossiuxea-de-campo}
\begin{description}[style=nextline,leftmargin=0.5cm,font=\bfseries]
\item[Coordenadas]  26°11'S 56°28'W (geocodificado via Nominatim).
\end{description}
\textbf{Fonte climática:} NASA POWER Climatology API (período 2001-2020)

\textbf{Fonte luminosa:} estimativa world atlas

\paragraph{Irradiação Solar
(kWh/\allowbreak m²/\allowbreak dia)}\label{irradiauxe7uxe3o-solar-kwhmuxb2dia}


\begin{center}
\resizebox{\linewidth}{!}{%
{\renewcommand{\arraystretch}{1.2}%
\begin{tabular}{@{}c|c|c|c|c|c|c|c|c|c|c|c|c@{}}
\toprule
Jan & Fev & Mar & Abr & Mai & Jun & Jul & Ago & Set & Out & Nov & Dez & Média \\
\midrule
6.71 & 6.20 & 5.39 & 4.34 & 3.26 & 2.77 & 3.10 & 3.79 & 4.42 & 5.34 &
6.39 & 6.76 & \textbf{4.87} \\
\bottomrule
\end{tabular}}%
}
\end{center}


\textbf{Inclinação solar recomendada:} 26° N (anual) · 36° N (inverno
jun-ago) · 16° N (verão nov-jan)

\paragraph{Precipitação (mm/\allowbreak mês)}\label{precipitauxe7uxe3o-mmmuxeas}


\begin{center}
\resizebox{\linewidth}{!}{%
{\renewcommand{\arraystretch}{1.2}%
\begin{tabular}{@{}c|c|c|c|c|c|c|c|c|c|c|c|c@{}}
\toprule
Jan & Fev & Mar & Abr & Mai & Jun & Jul & Ago & Set & Out & Nov & Dez & Total/\allowbreak ano \\
\midrule
135 & 129 & 133 & 159 & 166 & 94 & 78 & 61 & 93 & 189 & 188 & 164 &
\textbf{1590 mm} \\
\bottomrule
\end{tabular}}%
}
\end{center}


\paragraph{Poluição Luminosa}\label{poluiuxe7uxe3o-luminosa}

\begin{longtable}[]{@{}ll@{}}
\toprule\noalign{}
Parâmetro & Valor \\
\midrule\noalign{}
\endhead
\bottomrule\noalign{}
\endlastfoot
Escala Bortle & 3 --- Céu rural \\
Radiância artificial & 0.8 nW/\allowbreak cm²/\allowbreak sr \\
\end{longtable}
\paragraph{Solo (SoilGrids 2.0, média ponderada 0--30
cm)}

\begin{longtable}[]{@{}ll@{}}
\toprule\noalign{}
Parâmetro & Valor \\
\midrule\noalign{}
\endhead
\bottomrule\noalign{}
\endlastfoot
pH (H$_2$O) & 5.4 \\
Carbono orgânico (SOC) & 15.07 g/\allowbreak kg \\
Argila & 22.9 \% \\
Areia & 50.16 \% \\
Silte (calc.) & 26.9 \% \\
Densidade aparente & 1.4 g/\allowbreak cm³ \\
\textbf{Aptidão agrícola} & \textbf{Média} \\
\end{longtable}

\subsubsection{Indicadores Sociais}

\textbf{Fontes:} PNUD 2020, INE\footnote{\fineINE} Censo 2022, INE\footnote{\fineINE} EPHC 2023, Ministerio
Público 2024\\
\textbf{Nota:} Valores marcados como \emph{(dept.)} referem-se ao
departamento de Caazapá; valores \emph{(dist.)} são específicos deste
distrito. Dados marcados \emph{(est.)} são estimativas.

\begin{longtable}[]{@{}
  >{\raggedright\arraybackslash}p{(\linewidth - 4\tabcolsep) * \real{0.4231}}
  >{\raggedright\arraybackslash}p{(\linewidth - 4\tabcolsep) * \real{0.2692}}
  >{\raggedright\arraybackslash}p{(\linewidth - 4\tabcolsep) * \real{0.3077}}@{}}
\toprule\noalign{}
\begin{minipage}[b]{\linewidth}\raggedright
Indicador
\end{minipage} & \begin{minipage}[b]{\linewidth}\raggedright
Valor
\end{minipage} & \begin{minipage}[b]{\linewidth}\raggedright
Âmbito
\end{minipage} \\
\midrule\noalign{}
\endhead
\bottomrule\noalign{}
\endlastfoot
IDH (2020) & 0,648 & dept. \\
Ranking IDH nacional & 18/\allowbreak 18 & dept. \\
Esperança de vida & 69,5 anos & dept. \\
Escolaridade média & 8.0 anos & dist. \\
RNB per capita (USD PPA) & 5.200 & dept. \\
Pobreza monetária (\%) & 34,7\% & dept. \\
Pobreza extrema (\%) & 11,3\% & dept. \\
Índice de Gini & 0,472 & dept. \\
Acesso a água potável (\%) & 60,1\% & dept. \\
Acesso a saneamento (\%) & 57,4\% & dept. \\
Taxa de homicídios (est., /\allowbreak 100k hab) & \textasciitilde4--6 (estimativa,
próxima à média nacional) & dept. \\
Índice de segurança & médio & dept. \\
População (Censo 2022) & 3.461 hab. & dist. \\
Idade mediana & 34.0 anos & dist. \\
Taxa de fecundidade & N/\allowbreak D & dist. \\
\end{longtable}
\subsubsection{Combustível}

\textbf{Referência:} PETROPAR /\allowbreak  postos locais (2024) \textbf{Tipo de
localidade:} Interior

\begin{longtable}[]{@{}lll@{}}
\toprule\noalign{}
Combustível & USD/\allowbreak litro & Gs/\allowbreak litro (aprox.) \\
\midrule\noalign{}
\endhead
\bottomrule\noalign{}
\endlastfoot
Gasolina 93 oct & 0.99 & 7,326 \\
Gasolina 97 oct (premium) & 1.09 & 8,066 \\
Diesel & 0.91 & 6,734 \\
\end{longtable}

\begin{quote}
Preços podem variar ±5\% conforme posto e sazonalidade. Chaco e interior
remoto apresentam maior variação.
\end{quote}

\subsubsection{Cobertura Celular}

\textbf{Fonte:} CONATEL PY /\allowbreak  operadoras (2024)

\begin{longtable}[]{@{}ll@{}}
\toprule\noalign{}
Parâmetro & Valor \\
\midrule\noalign{}
\endhead
\bottomrule\noalign{}
\endlastfoot
Cobertura 4G população (dept.) & 78\% \\
Cobertura 4G área rural & 55\% \\
Melhor operadora & Tigo \\
Qualidade rural & moderada \\
\end{longtable}

\begin{quote}
Para áreas rurais fora do núcleo urbano, recomenda-se chip Tigo como
principal e Personal como backup.
\end{quote}

\subsubsection{Internet}

\textbf{Fonte:} CONATEL /\allowbreak  Speedtest Ookla (2024)

\begin{longtable}[]{@{}ll@{}}
\toprule\noalign{}
Parâmetro & Valor \\
\midrule\noalign{}
\endhead
\bottomrule\noalign{}
\endlastfoot
Velocidade média download & 32 Mbps \\
Domicílios com internet (dept.) & 42\% \\
Tecnologia predominante & rádio \\
Opção rural & Starlink disponível (\textasciitilde USD 44/\allowbreak mês) \\
\end{longtable}

\subsubsection{Mercado Imobiliário e Terra
Rural}

\textbf{Fonte:} INDERT /\allowbreak  Clasificados.com.py (2024)

\begin{longtable}[]{@{}ll@{}}
\toprule\noalign{}
Tipo & Referência \\
\midrule\noalign{}
\endhead
\bottomrule\noalign{}
\endlastfoot
Terra agrícola alta prod. (USD/\allowbreak ha) & 5,500 \\
Imóvel urbano (USD/\allowbreak m²) & 550 \\
Aluguel 2 quartos (USD/\allowbreak mês) & 220 \\
\end{longtable}

\begin{quote}
Valores de referência departamental. Localidades menores podem ter
preços 20--40\% abaixo da capital departamental. \#\#\# Saúde
\end{quote}

\textbf{Fonte:} MSPBS /\allowbreak  IPS Paraguay (2024-2026), consolidação
departamental e proxy local

\begin{longtable}[]{@{}ll@{}}
\toprule\noalign{}
Serviço & Disponibilidade \\
\midrule\noalign{}
\endhead
\bottomrule\noalign{}
\endlastfoot
USF /\allowbreak  Posto de Saúde & sim \\
Hospital Regional & não \\
IPS (seguro social) & não \\
Farmácia & sim \\
Distância ao hospital de referência & \textasciitilde13 km (Caazapa) \\
\end{longtable}

\textbf{Principais estabelecimentos:} USF local; referencia hospitalar
em Caazapa

\textbf{Observação para imigrantes:} Atendimento primario local ou em
raio curto; casos de maior complexidade seguem para Caazapa. Cobertura
privada continua recomendada para especialidades e urgências de maior
complexidade.
\newpage
\secaoDiagnostico{San Juan Nepomuceno}{\mapaDistrito{0607}}{San Juan Nepomuceno é o motor econômico de Caazapá. Oferece a melhor
infraestrutura de serviços do departamento, mas sua escala urbana requer
maior planejamento de segurança e logística em crises do que distritos
puramente rurais.

\subsubsection{Recomendação
Técnica}

Ideal para quem busca uma base operacional com acesso a serviços e mão
de obra, sem abrir mão da segurança rural.

\subsubsection{Analise de Sensibilidade}

\begin{itemize}
\tightlist
\item
  Censo 2022 INE\footnote{\fineINE}.
\item
  Relatórios de infraestrutura MOPC.
\item
  Dados de segurança do Ministério do Interior.
\end{itemize}}
\begin{table}[ht!]
\small \begin{tabular}{p{3.0cm}p{0.8cm}p{6.0cm}} \toprule
\textbf{Critério} & \textbf{Nota} & \textbf{Justificativa} \\ \midrule
A - Ameaças Estratégicas & 4.2 & - \\
B - Risco Social & 6.0 & - \\
C - Riscos Naturais & 7.0 & - \\
D - Autossuficiência & 6.0 & - \\
E - Institucional & 6.0 & - \\
\midrule
 \textbf{GSS FINAL} & \textbf{5.7} & \textbf{Seguro (Potencial moderado de resiliência).} \\ \bottomrule
 \end{tabular}
\end{table}
\subsubsection{Vulnerabilidades e
Mitigação}\label{vuln-San_Juan_Nepomuceno-vulnerabilidades-e-mitigauxe7uxe3o}

\begin{itemize}
\item
  \textbf{Pressão Urbana:} Em cenários de crise, pode atrair refugiados
  de áreas rurais vizinhas. Mitigação: Fortalecimento da segurança
  comunitária e estoques reguladores.
\item
  \textbf{Logística:} Dependência de poucas rotas pavimentadas para
  escoamento. Mitigação: Uso de aeródromo local para emergências.
\item
  INE\footnote{\fineINE} Censo 2022: 28.233 habitantes.
\item
  Preço da Terra: US\$ 5.000 a US\$ 6.500 (Terras agrícolas
  mecanizadas).
\item
  Homicídios: 6,5/\allowbreak 100k (Regional).
\end{itemize}
\subsubsection{Dossiê de Campo}\label{dossie-San_Juan_Nepomuceno-dossiuxea-de-campo}
\begin{description}[style=nextline,leftmargin=0.5cm,font=\bfseries]
\item[Coordenadas]  26°06'S 55°58'W.
\item[Topografia]  Relevo ondulado com presença de vales férteis. Área de \textasciitilde 1.011 km².
\item[Alvos Estrategicos]  Principal centro comercial e populacional do departamento. Localizado na interseção de rotas que ligam ao leste e sul. Sem alvos militares diretos.
\item[Fallout]  Ventos predominantes do Norte/\allowbreak Nordeste.
\item[Populacao]  28.233 (Censo 2022 Final). Cidade mais populosa de Caazapá.
\item[Densidade]  \textasciitilde 27,9 hab/\allowbreak km².
\item[Vias de Acesso]  Acesso pela Ruta PY08 e conexões asfaltadas recentes para o norte.
\item[Servicos]  Hospital Distrital e centros educacionais consolidados. Comércio diversificado.
\item[Hidrologia]  Diversos arroios cruzam o distrito. Risco moderado de interrupção de pontes rurais em eventos de precipitação severa.
\item[Sismicidade]  Baixa.
\item[Solo]  Ultisols e Oxisols (Terra Roxa no leste), alta fertilidade.
\item[Agua]  Sistemas de Juntas de Saneamiento e poços profundos.
\item[Energia]  Rede ANDE\footnote{\fineANDE}.
\item[Producao]  Soja, milho, cana-de-açúcar e pecuária.
\item[Seguranca]  Taxa de homicídios regional estimada em 6,5/\allowbreak 100k hab. Segurança rural estável, com foco em conflitos de vizinhança.
\item[Clima Social]  Sociedade vibrante com forte identidade agrícola e comercial.
\end{description}
\textbf{Fonte climática:} NASA POWER Climatology API (período 2001-2020)

\textbf{Fonte luminosa:} estimativa world atlas

\paragraph{Irradiação Solar
(kWh/\allowbreak m²/\allowbreak dia)}\label{irradiauxe7uxe3o-solar-kwhmuxb2dia}


\begin{center}
\resizebox{\linewidth}{!}{%
{\renewcommand{\arraystretch}{1.2}%
\begin{tabular}{@{}c|c|c|c|c|c|c|c|c|c|c|c|c@{}}
\toprule
Jan & Fev & Mar & Abr & Mai & Jun & Jul & Ago & Set & Out & Nov & Dez & Média \\
\midrule
6.56 & 6.06 & 5.32 & 4.33 & 3.21 & 2.76 & 3.08 & 3.81 & 4.39 & 5.25 &
6.31 & 6.67 & \textbf{4.81} \\
\bottomrule
\end{tabular}}%
}
\end{center}


\textbf{Inclinação solar recomendada:} 26° N (anual) · 36° N (inverno
jun-ago) · 16° N (verão nov-jan)

\paragraph{Precipitação (mm/\allowbreak mês)}\label{precipitauxe7uxe3o-mmmuxeas}


\begin{center}
\resizebox{\linewidth}{!}{%
{\renewcommand{\arraystretch}{1.2}%
\begin{tabular}{@{}c|c|c|c|c|c|c|c|c|c|c|c|c@{}}
\toprule
Jan & Fev & Mar & Abr & Mai & Jun & Jul & Ago & Set & Out & Nov & Dez & Total/\allowbreak ano \\
\midrule
142 & 127 & 131 & 161 & 166 & 102 & 83 & 64 & 102 & 193 & 184 & 169 &
\textbf{1625 mm} \\
\bottomrule
\end{tabular}}%
}
\end{center}


\paragraph{Poluição Luminosa}\label{poluiuxe7uxe3o-luminosa}

\begin{longtable}[]{@{}ll@{}}
\toprule\noalign{}
Parâmetro & Valor \\
\midrule\noalign{}
\endhead
\bottomrule\noalign{}
\endlastfoot
Escala Bortle & 3 --- Céu rural \\
Radiância artificial & 0.8 nW/\allowbreak cm²/\allowbreak sr \\
\end{longtable}
\paragraph{Solo (SoilGrids 2.0, média ponderada 0--30
cm)}

\begin{longtable}[]{@{}ll@{}}
\toprule\noalign{}
Parâmetro & Valor \\
\midrule\noalign{}
\endhead
\bottomrule\noalign{}
\endlastfoot
pH (H$_2$O) & 5.2 \\
Carbono orgânico (SOC) & 18.35 g/\allowbreak kg \\
Argila & 31.1 \% \\
Areia & 37.77 \% \\
Silte (calc.) & 31.1 \% \\
Densidade aparente & 1.32 g/\allowbreak cm³ \\
\textbf{Aptidão agrícola} & \textbf{Média} \\
\end{longtable}

\begin{itemize}
\tightlist
\item
  \textbf{Agua:} Sistemas de Juntas de Saneamiento e poços profundos.
\item
  \textbf{Energia:} Rede ANDE\footnote{\fineANDE}.
\item
  \textbf{Producao:} Soja, milho, cana-de-açúcar e pecuária.
\end{itemize}
\subsubsection{Indicadores Sociais}

\textbf{Fontes:} PNUD 2020, INE\footnote{\fineINE} Censo 2022, INE\footnote{\fineINE} EPHC 2023, Ministerio
Público 2024\\
\textbf{Nota:} Valores marcados como \emph{(dept.)} referem-se ao
departamento de Caazapá; valores \emph{(dist.)} são específicos deste
distrito. Dados marcados \emph{(est.)} são estimativas.

\begin{longtable}[]{@{}
  >{\raggedright\arraybackslash}p{(\linewidth - 4\tabcolsep) * \real{0.4231}}
  >{\raggedright\arraybackslash}p{(\linewidth - 4\tabcolsep) * \real{0.2692}}
  >{\raggedright\arraybackslash}p{(\linewidth - 4\tabcolsep) * \real{0.3077}}@{}}
\toprule\noalign{}
\begin{minipage}[b]{\linewidth}\raggedright
Indicador
\end{minipage} & \begin{minipage}[b]{\linewidth}\raggedright
Valor
\end{minipage} & \begin{minipage}[b]{\linewidth}\raggedright
Âmbito
\end{minipage} \\
\midrule\noalign{}
\endhead
\bottomrule\noalign{}
\endlastfoot
IDH (2020) & 0,648 & dept. \\
Ranking IDH nacional & 18/\allowbreak 18 & dept. \\
Esperança de vida & 69,5 anos & dept. \\
Escolaridade média & 8.4 anos & dist. \\
RNB per capita (USD PPA) & 5.200 & dept. \\
Pobreza monetária (\%) & 34,7\% & dept. \\
Pobreza extrema (\%) & 11,3\% & dept. \\
Índice de Gini & 0,472 & dept. \\
Acesso a água potável (\%) & 60,1\% & dept. \\
Acesso a saneamento (\%) & 57,4\% & dept. \\
Taxa de homicídios (est., /\allowbreak 100k hab) & \textasciitilde4--6 (estimativa,
próxima à média nacional) & dept. \\
Índice de segurança & médio & dept. \\
População (Censo 2022) & 28.233 hab. & dist. \\
Idade mediana & 28.0 anos & dist. \\
Taxa de fecundidade & N/\allowbreak D & dist. \\
\end{longtable}
\subsubsection{Combustível}

\textbf{Referência:} PETROPAR /\allowbreak  postos locais (2024) \textbf{Tipo de
localidade:} Interior

\begin{longtable}[]{@{}lll@{}}
\toprule\noalign{}
Combustível & USD/\allowbreak litro & Gs/\allowbreak litro (aprox.) \\
\midrule\noalign{}
\endhead
\bottomrule\noalign{}
\endlastfoot
Gasolina 93 oct & 0.99 & 7,326 \\
Gasolina 97 oct (premium) & 1.09 & 8,066 \\
Diesel & 0.91 & 6,734 \\
\end{longtable}

\begin{quote}
Preços podem variar ±5\% conforme posto e sazonalidade. Chaco e interior
remoto apresentam maior variação.
\end{quote}

\subsubsection{Cobertura Celular}

\textbf{Fonte:} CONATEL PY /\allowbreak  operadoras (2024)

\begin{longtable}[]{@{}ll@{}}
\toprule\noalign{}
Parâmetro & Valor \\
\midrule\noalign{}
\endhead
\bottomrule\noalign{}
\endlastfoot
Cobertura 4G população (dept.) & 78\% \\
Cobertura 4G área rural & 55\% \\
Melhor operadora & Tigo \\
Qualidade rural & moderada \\
\end{longtable}

\begin{quote}
Para áreas rurais fora do núcleo urbano, recomenda-se chip Tigo como
principal e Personal como backup.
\end{quote}

\subsubsection{Internet}

\textbf{Fonte:} CONATEL /\allowbreak  Speedtest Ookla (2024)

\begin{longtable}[]{@{}ll@{}}
\toprule\noalign{}
Parâmetro & Valor \\
\midrule\noalign{}
\endhead
\bottomrule\noalign{}
\endlastfoot
Velocidade média download & 32 Mbps \\
Domicílios com internet (dept.) & 42\% \\
Tecnologia predominante & rádio \\
Opção rural & Starlink disponível (\textasciitilde USD 44/\allowbreak mês) \\
\end{longtable}

\subsubsection{Mercado Imobiliário e Terra
Rural}

\textbf{Fonte:} INDERT /\allowbreak  Clasificados.com.py (2024)

\begin{longtable}[]{@{}ll@{}}
\toprule\noalign{}
Tipo & Referência \\
\midrule\noalign{}
\endhead
\bottomrule\noalign{}
\endlastfoot
Terra agrícola alta prod. (USD/\allowbreak ha) & 5,500 \\
Imóvel urbano (USD/\allowbreak m²) & 550 \\
Aluguel 2 quartos (USD/\allowbreak mês) & 220 \\
\end{longtable}

\begin{quote}
Valores de referência departamental. Localidades menores podem ter
preços 20--40\% abaixo da capital departamental. \#\#\# Saúde
\end{quote}

\textbf{Fonte:} MSPBS /\allowbreak  IPS Paraguay (2024-2026), consolidação
departamental e proxy local

\begin{longtable}[]{@{}ll@{}}
\toprule\noalign{}
Serviço & Disponibilidade \\
\midrule\noalign{}
\endhead
\bottomrule\noalign{}
\endlastfoot
USF /\allowbreak  Posto de Saúde & sim \\
Hospital Regional & sim \\
IPS (seguro social) & não \\
Farmácia & sim \\
Distância ao hospital de referência & \textasciitilde55 km (Caazapa) \\
\end{longtable}

\textbf{Principais estabelecimentos:} USF local; referencia hospitalar
em Caazapa

\textbf{Observação para imigrantes:} Atendimento primario local ou em
raio curto; casos de maior complexidade seguem para Caazapa. Cobertura
privada continua recomendada para especialidades e urgências de maior
complexidade.
\newpage
\secaoDiagnostico{Tavarai}{\mapaDistrito{0409}}{Tavarai (Caazapa) apresenta perfil operacional consolidado para
comparacao territorial, com leitura conjunta de infraestrutura, risco
hidroclimatico e capacidade de continuidade de servicos essenciais.

\subsubsection{Pre-condicoes Minimas de
Decisao}

\begin{itemize}
\tightlist
\item
  Atualizar indicadores criticos em ciclo trimestral.
\item
  Confirmar trafegabilidade e contingencia hidrica/\allowbreak energetica local.
\item
  Validar checklist de resiliencia domiciliar/\allowbreak logistica antes de decisao
  final.
\end{itemize}

\subsubsection{Recomendação
Técnica}

Uso recomendado como candidato em analise multicriterio, condicionado ao
cumprimento das pre-condicoes acima. Referência atual de score: GSS 8.3;
risco predominante: baixo.

}
\begin{table}[ht!]
\small \begin{tabular}{p{3.0cm}p{0.8cm}p{6.0cm}} \toprule
\textbf{Critério} & \textbf{Nota} & \textbf{Justificativa} \\ \midrule
A - Ameaças Estratégicas & 9.1 & - \\
B - Risco Social & 8.1 & - \\
C - Riscos Naturais & 7.9 & - \\
D - Autossuficiência & 7.9 & - \\
E - Institucional & 8.3 & - \\
\midrule
 \textbf{GSS FINAL} & \textbf{8.3} & \textbf{Seguro (bom potencial com preparacao).} \\ \bottomrule
 \end{tabular}
\end{table}
\subsubsection{Vulnerabilidades e
Mitigação}\label{vuln-Tavarai-vulnerabilidades-e-mitigauxe7uxe3o}

\begin{itemize}
\tightlist
\item
  Exposicao hidrometeorologica controlada (\textless45/\allowbreak 100).
\item
  Conectividade viaria favoravel (\textgreater=70/\allowbreak 100).
\item
  Estoque de contingencia para 14 dias (agua/\allowbreak alimentos/\allowbreak combustivel),
  priorizado quando risco hidro=26/\allowbreak 100.
\item
  Duplicar rotas/\allowbreak logistica quando conectividade viaria=95/\allowbreak 100 for
  \textless70; manter rota secundaria mapeada e testada trimestralmente.
\item
  Plano de energia resiliente com autonomia minima de 72h
  (geracao/\allowbreak backup), revisao semestral.
\end{itemize}

\subsubsection{Dossiê de Campo}\label{dossie-Tavarai-dossiuxea-de-campo}
\begin{description}[style=nextline,leftmargin=0.5cm,font=\bfseries]
\item[Coordenadas]  26°11'S 56°34'W (geocodificado via Nominatim).
\end{description}
\textbf{Fonte climática:} NASA POWER Climatology API (período 2001-2020)

\textbf{Fonte luminosa:} estimativa world atlas

\paragraph{Irradiação Solar
(kWh/\allowbreak m²/\allowbreak dia)}\label{irradiauxe7uxe3o-solar-kwhmuxb2dia}


\begin{center}
\resizebox{\linewidth}{!}{%
{\renewcommand{\arraystretch}{1.2}%
\begin{tabular}{@{}c|c|c|c|c|c|c|c|c|c|c|c|c@{}}
\toprule
Jan & Fev & Mar & Abr & Mai & Jun & Jul & Ago & Set & Out & Nov & Dez & Média \\
\midrule
6.71 & 6.20 & 5.39 & 4.34 & 3.26 & 2.77 & 3.10 & 3.79 & 4.42 & 5.34 &
6.39 & 6.76 & \textbf{4.87} \\
\bottomrule
\end{tabular}}%
}
\end{center}


\textbf{Inclinação solar recomendada:} 26° N (anual) · 36° N (inverno
jun-ago) · 16° N (verão nov-jan)

\paragraph{Precipitação (mm/\allowbreak mês)}\label{precipitauxe7uxe3o-mmmuxeas}


\begin{center}
\resizebox{\linewidth}{!}{%
{\renewcommand{\arraystretch}{1.2}%
\begin{tabular}{@{}c|c|c|c|c|c|c|c|c|c|c|c|c@{}}
\toprule
Jan & Fev & Mar & Abr & Mai & Jun & Jul & Ago & Set & Out & Nov & Dez & Total/\allowbreak ano \\
\midrule
135 & 129 & 133 & 159 & 166 & 94 & 78 & 61 & 93 & 189 & 188 & 164 &
\textbf{1590 mm} \\
\bottomrule
\end{tabular}}%
}
\end{center}


\paragraph{Poluição Luminosa}\label{poluiuxe7uxe3o-luminosa}

\begin{longtable}[]{@{}ll@{}}
\toprule\noalign{}
Parâmetro & Valor \\
\midrule\noalign{}
\endhead
\bottomrule\noalign{}
\endlastfoot
Escala Bortle & 3 --- Céu rural \\
Radiância artificial & 0.8 nW/\allowbreak cm²/\allowbreak sr \\
\end{longtable}
\paragraph{Solo (SoilGrids 2.0, média ponderada 0--30
cm)}

\begin{longtable}[]{@{}ll@{}}
\toprule\noalign{}
Parâmetro & Valor \\
\midrule\noalign{}
\endhead
\bottomrule\noalign{}
\endlastfoot
pH (H$_2$O) & 5.3 \\
Carbono orgânico (SOC) & 19.12 g/\allowbreak kg \\
Argila & 24.55 \% \\
Areia & 48.54 \% \\
Silte (calc.) & 26.9 \% \\
Densidade aparente & 1.35 g/\allowbreak cm³ \\
\textbf{Aptidão agrícola} & \textbf{Média} \\
\end{longtable}

\subsubsection{Indicadores Sociais}

\textbf{Fontes:} PNUD 2020, INE\footnote{\fineINE} Censo 2022, INE\footnote{\fineINE} EPHC 2023, Ministerio
Público 2024\\
\textbf{Nota:} Valores marcados como \emph{(dept.)} referem-se ao
departamento de Caazapá; valores \emph{(dist.)} são específicos deste
distrito. Dados marcados \emph{(est.)} são estimativas.

\begin{longtable}[]{@{}
  >{\raggedright\arraybackslash}p{(\linewidth - 4\tabcolsep) * \real{0.4231}}
  >{\raggedright\arraybackslash}p{(\linewidth - 4\tabcolsep) * \real{0.2692}}
  >{\raggedright\arraybackslash}p{(\linewidth - 4\tabcolsep) * \real{0.3077}}@{}}
\toprule\noalign{}
\begin{minipage}[b]{\linewidth}\raggedright
Indicador
\end{minipage} & \begin{minipage}[b]{\linewidth}\raggedright
Valor
\end{minipage} & \begin{minipage}[b]{\linewidth}\raggedright
Âmbito
\end{minipage} \\
\midrule\noalign{}
\endhead
\bottomrule\noalign{}
\endlastfoot
IDH (2020) & 0,648 & dept. \\
Ranking IDH nacional & 18/\allowbreak 18 & dept. \\
Esperança de vida & 69,5 anos & dept. \\
Escolaridade média & 6,4 anos & dept. \\
RNB per capita (USD PPA) & 5.200 & dept. \\
Pobreza monetária (\%) & 34,7\% & dept. \\
Pobreza extrema (\%) & 11,3\% & dept. \\
Índice de Gini & 0,472 & dept. \\
Acesso a água potável (\%) & 60,1\% & dept. \\
Acesso a saneamento (\%) & 57,4\% & dept. \\
Taxa de homicídios (est., /\allowbreak 100k hab) & \textasciitilde4--6 (estimativa,
próxima à média nacional) & dept. \\
Índice de segurança & médio & dept. \\
População (Censo 2022) & N/\allowbreak D & dist. \\
Idade mediana & N/\allowbreak D & dist. \\
Taxa de fecundidade & N/\allowbreak D & dist. \\
\end{longtable}
\subsubsection{Combustível}

\textbf{Referência:} PETROPAR /\allowbreak  postos locais (2024) \textbf{Tipo de
localidade:} Interior

\begin{longtable}[]{@{}lll@{}}
\toprule\noalign{}
Combustível & USD/\allowbreak litro & Gs/\allowbreak litro (aprox.) \\
\midrule\noalign{}
\endhead
\bottomrule\noalign{}
\endlastfoot
Gasolina 93 oct & 0.99 & 7,326 \\
Gasolina 97 oct (premium) & 1.09 & 8,066 \\
Diesel & 0.91 & 6,734 \\
\end{longtable}

\begin{quote}
Preços podem variar ±5\% conforme posto e sazonalidade. Chaco e interior
remoto apresentam maior variação.
\end{quote}

\subsubsection{Cobertura Celular}

\textbf{Fonte:} CONATEL PY /\allowbreak  operadoras (2024)

\begin{longtable}[]{@{}ll@{}}
\toprule\noalign{}
Parâmetro & Valor \\
\midrule\noalign{}
\endhead
\bottomrule\noalign{}
\endlastfoot
Cobertura 4G população (dept.) & 78\% \\
Cobertura 4G área rural & 55\% \\
Melhor operadora & Tigo \\
Qualidade rural & moderada \\
\end{longtable}

\begin{quote}
Para áreas rurais fora do núcleo urbano, recomenda-se chip Tigo como
principal e Personal como backup.
\end{quote}

\subsubsection{Internet}

\textbf{Fonte:} CONATEL /\allowbreak  Speedtest Ookla (2024)

\begin{longtable}[]{@{}ll@{}}
\toprule\noalign{}
Parâmetro & Valor \\
\midrule\noalign{}
\endhead
\bottomrule\noalign{}
\endlastfoot
Velocidade média download & 32 Mbps \\
Domicílios com internet (dept.) & 42\% \\
Tecnologia predominante & rádio \\
Opção rural & Starlink disponível (\textasciitilde USD 44/\allowbreak mês) \\
\end{longtable}

\subsubsection{Mercado Imobiliário e Terra
Rural}

\textbf{Fonte:} INDERT /\allowbreak  Clasificados.com.py (2024)

\begin{longtable}[]{@{}ll@{}}
\toprule\noalign{}
Tipo & Referência \\
\midrule\noalign{}
\endhead
\bottomrule\noalign{}
\endlastfoot
Terra agrícola alta prod. (USD/\allowbreak ha) & 5,500 \\
Imóvel urbano (USD/\allowbreak m²) & 550 \\
Aluguel 2 quartos (USD/\allowbreak mês) & 220 \\
\end{longtable}

\begin{quote}
Valores de referência departamental. Localidades menores podem ter
preços 20--40\% abaixo da capital departamental. \#\#\# Saúde
\end{quote}

\textbf{Fonte:} MSPBS /\allowbreak  IPS Paraguay (2024-2026), consolidação
departamental e proxy local

\begin{longtable}[]{@{}ll@{}}
\toprule\noalign{}
Serviço & Disponibilidade \\
\midrule\noalign{}
\endhead
\bottomrule\noalign{}
\endlastfoot
USF /\allowbreak  Posto de Saúde & sim \\
Hospital Regional & não \\
IPS (seguro social) & não \\
Farmácia & sim \\
Distância ao hospital de referência & \textasciitilde24 km (Caazapa) \\
\end{longtable}

\textbf{Principais estabelecimentos:} USF local; referencia hospitalar
em Caazapa

\textbf{Observação para imigrantes:} Atendimento primario local ou em
raio curto; casos de maior complexidade seguem para Caazapa. Cobertura
privada continua recomendada para especialidades e urgências de maior
complexidade.
\newpage
\secaoDiagnostico{Yuty}{\mapaDistrito{0610}}{Yuty é uma localidade de altíssima resiliência devido à sua baixa
densidade populacional e grande capacidade produtiva, especialmente em
pecuária e arroz. É um local discreto, seguro e com boa infraestrutura
básica, ideal para retiros estratégicos.

\subsubsection{Recomendação
Técnica}

Altamente recomendado para quem busca isolamento com segurança jurídica
e social.

\subsubsection{Analise de Sensibilidade}

\begin{itemize}
\tightlist
\item
  Dados demográficos validados INE\footnote{\fineINE} 2022.
\item
  Mapeamento hidrológico (DMH/\allowbreak SEN\footnote{\fineSEN}).
\item
  Relatórios de uso da terra (MAG).
\end{itemize}}
\begin{table}[ht!]
\small \begin{tabular}{p{3.0cm}p{0.8cm}p{6.0cm}} \toprule
\textbf{Critério} & \textbf{Nota} & \textbf{Justificativa} \\ \midrule
A - Ameaças Estratégicas & 6.3 & - \\
B - Risco Social & 6.5 & - \\
C - Riscos Naturais & 7.5 & - \\
D - Autossuficiência & 7.5 & - \\
E - Institucional & 7.0 & - \\
\midrule
 \textbf{GSS FINAL} & \textbf{6.9} & \textbf{Seguro (Potencial elevado de resiliência).} \\ \bottomrule
 \end{tabular}
\end{table}
\subsubsection{Vulnerabilidades e
Mitigação}\label{vuln-Yuty-vulnerabilidades-e-mitigauxe7uxe3o}

\begin{itemize}
\item
  \textbf{Alagamentos:} Áreas de pastagem baixas próximas ao Rio
  Tebicuary. Mitigação: Mapeamento de zonas de inundação e construção em
  cotas elevadas.
\item
  \textbf{Acesso:} Rodovias de terra podem ficar intrafegáveis.
  Mitigação: Veículos 4x4 e autonomia logística de 14 dias.
\item
  INE\footnote{\fineINE} Censo 2022: 14.924 habitantes.
\item
  Preço da Terra: US\$ 2.500 a US\$ 3.500 (Pecuária/\allowbreak Campos).
\item
  Homicídios: 6,5/\allowbreak 100k (Regional).
\end{itemize}
\subsubsection{Dossiê de Campo}\label{dossie-Yuty-dossiuxea-de-campo}
\begin{description}[style=nextline,leftmargin=0.5cm,font=\bfseries]
\item[Coordenadas]  26°36'S 56°15'W.
\item[Topografia]  Relevo de planície com suaves ondulações. Área de \textasciitilde 1.440 km².
\item[Alvos Estrategicos]  Localidade de baixa densidade no sul do departamento. Sem alvos militares ou infraestruturas críticas sensíveis.
\item[Fallout]  Ventos predominantes do Norte.
\item[Populacao]  14.924 (Censo 2022 Final).
\item[Densidade]  \textasciitilde 10,4 hab/\allowbreak km² (Excelente para isolamento e segurança social).
\item[Vias de Acesso]  Conexão pela Ruta PY08 e rotas departamentais. Presença histórica de linha ferroviária (atualmente inativa).
\item[Servicos]  Saúde básica e serviços municipais consolidados.
\item[Hidrologia]  Banhado pelo Rio Tebicuary e afluentes. Áreas baixas podem sofrer alagamentos, mas a zona urbana é elevada.
\item[Sismicidade]  Nula.
\item[Solo]  Arenas Cuarzosas e Ultisols. Excelente para pecuária e reflorestamento.
\item[Agua]  Recursos hídricos superficiais abundantes.
\item[Energia]  Rede ANDE\footnote{\fineANDE}.
\item[Producao]  Arroz, pecuária, e pequena agricultura de subsistência.
\item[Seguranca]  Uma das zonas mais pacíficas da Região Oriental. Baixíssima incidência de crimes violentos.
\item[Clima Social]  Estabilidade social tradicional. Comunidades rurais muito coesas.
\end{description}
\textbf{Fonte climática:} NASA POWER Climatology API (período 2001-2020)

\textbf{Fonte luminosa:} estimativa world atlas

\paragraph{Irradiação Solar
(kWh/\allowbreak m²/\allowbreak dia)}\label{irradiauxe7uxe3o-solar-kwhmuxb2dia}


\begin{center}
\resizebox{\linewidth}{!}{%
{\renewcommand{\arraystretch}{1.2}%
\begin{tabular}{@{}c|c|c|c|c|c|c|c|c|c|c|c|c@{}}
\toprule
Jan & Fev & Mar & Abr & Mai & Jun & Jul & Ago & Set & Out & Nov & Dez & Média \\
\midrule
6.71 & 6.20 & 5.39 & 4.34 & 3.26 & 2.77 & 3.10 & 3.79 & 4.42 & 5.34 &
6.39 & 6.76 & \textbf{4.87} \\
\bottomrule
\end{tabular}}%
}
\end{center}


\textbf{Inclinação solar recomendada:} 27° N (anual) · 37° N (inverno
jun-ago) · 17° N (verão nov-jan)

\paragraph{Precipitação (mm/\allowbreak mês)}\label{precipitauxe7uxe3o-mmmuxeas}


\begin{center}
\resizebox{\linewidth}{!}{%
{\renewcommand{\arraystretch}{1.2}%
\begin{tabular}{@{}c|c|c|c|c|c|c|c|c|c|c|c|c@{}}
\toprule
Jan & Fev & Mar & Abr & Mai & Jun & Jul & Ago & Set & Out & Nov & Dez & Total/\allowbreak ano \\
\midrule
141 & 123 & 136 & 171 & 156 & 100 & 82 & 71 & 100 & 194 & 188 & 167 &
\textbf{1629 mm} \\
\bottomrule
\end{tabular}}%
}
\end{center}


\paragraph{Poluição Luminosa}\label{poluiuxe7uxe3o-luminosa}

\begin{longtable}[]{@{}ll@{}}
\toprule\noalign{}
Parâmetro & Valor \\
\midrule\noalign{}
\endhead
\bottomrule\noalign{}
\endlastfoot
Escala Bortle & 3 --- Céu rural \\
Radiância artificial & 0.8 nW/\allowbreak cm²/\allowbreak sr \\
\end{longtable}
\paragraph{Solo (SoilGrids 2.0, média ponderada 0--30
cm)}

\begin{longtable}[]{@{}ll@{}}
\toprule\noalign{}
Parâmetro & Valor \\
\midrule\noalign{}
\endhead
\bottomrule\noalign{}
\endlastfoot
pH (H$_2$O) & 5.12 \\
Carbono orgânico (SOC) & 17.35 g/\allowbreak kg \\
Argila & 28.18 \% \\
Areia & 41.08 \% \\
Silte (calc.) & 30.7 \% \\
Densidade aparente & 1.38 g/\allowbreak cm³ \\
\textbf{Aptidão agrícola} & \textbf{Média} \\
\end{longtable}

\begin{itemize}
\tightlist
\item
  \textbf{Agua:} Recursos hídricos superficiais abundantes.
\item
  \textbf{Energia:} Rede ANDE\footnote{\fineANDE}.
\item
  \textbf{Producao:} Arroz, pecuária, e pequena agricultura de
  subsistência.
\end{itemize}
\subsubsection{Indicadores Sociais}

\textbf{Fontes:} PNUD 2020, INE\footnote{\fineINE} Censo 2022, INE\footnote{\fineINE} EPHC 2023, Ministerio
Público 2024\\
\textbf{Nota:} Valores marcados como \emph{(dept.)} referem-se ao
departamento de Caazapá; valores \emph{(dist.)} são específicos deste
distrito. Dados marcados \emph{(est.)} são estimativas.

\begin{longtable}[]{@{}
  >{\raggedright\arraybackslash}p{(\linewidth - 4\tabcolsep) * \real{0.4231}}
  >{\raggedright\arraybackslash}p{(\linewidth - 4\tabcolsep) * \real{0.2692}}
  >{\raggedright\arraybackslash}p{(\linewidth - 4\tabcolsep) * \real{0.3077}}@{}}
\toprule\noalign{}
\begin{minipage}[b]{\linewidth}\raggedright
Indicador
\end{minipage} & \begin{minipage}[b]{\linewidth}\raggedright
Valor
\end{minipage} & \begin{minipage}[b]{\linewidth}\raggedright
Âmbito
\end{minipage} \\
\midrule\noalign{}
\endhead
\bottomrule\noalign{}
\endlastfoot
IDH (2020) & 0,648 & dept. \\
Ranking IDH nacional & 18/\allowbreak 18 & dept. \\
Esperança de vida & 69,5 anos & dept. \\
Escolaridade média & 8.1 anos & dist. \\
RNB per capita (USD PPA) & 5.200 & dept. \\
Pobreza monetária (\%) & 34,7\% & dept. \\
Pobreza extrema (\%) & 11,3\% & dept. \\
Índice de Gini & 0,472 & dept. \\
Acesso a água potável (\%) & 60,1\% & dept. \\
Acesso a saneamento (\%) & 57,4\% & dept. \\
Taxa de homicídios (est., /\allowbreak 100k hab) & \textasciitilde4--6 (estimativa,
próxima à média nacional) & dept. \\
Índice de segurança & médio & dept. \\
População (Censo 2022) & 14.924 hab. & dist. \\
Idade mediana & 31.0 anos & dist. \\
Taxa de fecundidade & N/\allowbreak D & dist. \\
\end{longtable}
\subsubsection{Combustível}

\textbf{Referência:} PETROPAR /\allowbreak  postos locais (2024) \textbf{Tipo de
localidade:} Interior

\begin{longtable}[]{@{}lll@{}}
\toprule\noalign{}
Combustível & USD/\allowbreak litro & Gs/\allowbreak litro (aprox.) \\
\midrule\noalign{}
\endhead
\bottomrule\noalign{}
\endlastfoot
Gasolina 93 oct & 0.99 & 7,326 \\
Gasolina 97 oct (premium) & 1.09 & 8,066 \\
Diesel & 0.91 & 6,734 \\
\end{longtable}

\begin{quote}
Preços podem variar ±5\% conforme posto e sazonalidade. Chaco e interior
remoto apresentam maior variação.
\end{quote}

\subsubsection{Cobertura Celular}

\textbf{Fonte:} CONATEL PY /\allowbreak  operadoras (2024)

\begin{longtable}[]{@{}ll@{}}
\toprule\noalign{}
Parâmetro & Valor \\
\midrule\noalign{}
\endhead
\bottomrule\noalign{}
\endlastfoot
Cobertura 4G população (dept.) & 78\% \\
Cobertura 4G área rural & 55\% \\
Melhor operadora & Tigo \\
Qualidade rural & moderada \\
\end{longtable}

\begin{quote}
Para áreas rurais fora do núcleo urbano, recomenda-se chip Tigo como
principal e Personal como backup.
\end{quote}

\subsubsection{Internet}

\textbf{Fonte:} CONATEL /\allowbreak  Speedtest Ookla (2024)

\begin{longtable}[]{@{}ll@{}}
\toprule\noalign{}
Parâmetro & Valor \\
\midrule\noalign{}
\endhead
\bottomrule\noalign{}
\endlastfoot
Velocidade média download & 32 Mbps \\
Domicílios com internet (dept.) & 42\% \\
Tecnologia predominante & rádio \\
Opção rural & Starlink disponível (\textasciitilde USD 44/\allowbreak mês) \\
\end{longtable}

\subsubsection{Mercado Imobiliário e Terra
Rural}

\textbf{Fonte:} INDERT /\allowbreak  Clasificados.com.py (2024)

\begin{longtable}[]{@{}ll@{}}
\toprule\noalign{}
Tipo & Referência \\
\midrule\noalign{}
\endhead
\bottomrule\noalign{}
\endlastfoot
Terra agrícola alta prod. (USD/\allowbreak ha) & 5,500 \\
Imóvel urbano (USD/\allowbreak m²) & 550 \\
Aluguel 2 quartos (USD/\allowbreak mês) & 220 \\
\end{longtable}

\begin{quote}
Valores de referência departamental. Localidades menores podem ter
preços 20--40\% abaixo da capital departamental. \#\#\# Saúde
\end{quote}

\textbf{Fonte:} MSPBS /\allowbreak  IPS Paraguay (2024-2026), consolidação
departamental e proxy local

\begin{longtable}[]{@{}ll@{}}
\toprule\noalign{}
Serviço & Disponibilidade \\
\midrule\noalign{}
\endhead
\bottomrule\noalign{}
\endlastfoot
USF /\allowbreak  Posto de Saúde & sim \\
Hospital Regional & não \\
IPS (seguro social) & não \\
Farmácia & sim \\
Distância ao hospital de referência & \textasciitilde60 km (Caazapa) \\
\end{longtable}

\textbf{Principais estabelecimentos:} USF local; referencia hospitalar
em Caazapa

\textbf{Observação para imigrantes:} Atendimento primario local ou em
raio curto; casos de maior complexidade seguem para Caazapa. Cobertura
privada continua recomendada para especialidades e urgências de maior
complexidade.
